% =====================================================================
%  R. Nielsen, TO SLAGS UNDERVISNING, I-II
%  In: For Ide og Virkelighed. Et Tidsskrift, udgivet af R. Nielsen, Bj. Bjørnson og Rud. Schmidt.  Aargang 1869.
%  Kjøbenhavn: Chr. Steen & Søns Forlag; Louis Kleins Bogtrykkeri.  pp. 297-323, pp. 409-441.  Antiqua.
%
%  TRANSCRIPTION -- GENERATED; DO NOT EDIT BY HAND.  The five Nielsen essays of the 1869 volume are built by ONE pipeline
%  (ebook-method/for-idee-1869/) into a combined working file, which split5.py cuts into the five transcriptions.
%  Rebuild from the repo root:   sh ebook-method/for-idee-1869/rebuild.sh
%
%  STATUS: COMPLETE (2026-10-13): the volume's title page and the whole essay, pp. 297-323, pp. 409-441, with its notes.
%
%  EDITION.  The only printing.  Text of record: the Internet Archive's scan of the University of Toronto copy
%  (forideogvirkelig00kjobuoft), ~/bibliotek/Nielsen, Rasmus/1869-for-idee-og-virkelighed.pdf (SCANS.tsv), 886 PDF pages,
%  400 ppi colour; text layer = the IA's ABBYY FineReader OCR.  Printed page N = PDF N + 10 to p. 688; a plate is bound in
%  before p. 689, so pp. 689-706 = PDF N + 12.  Title page PDF 7.  No e-book, no other scan, no earlier transcription found
%  (checked 2026-10-13: the repo, the web, the Internet Archive; HathiTrust and Google Books could not be queried that day).
%
%  ACCURACY (2026-10-13; one collation for all five essays, 189 pp., ~48,000 words).  The words are those of the ABBYY layer
%  (draft.py), collated word by word with an independent tesseract reading of the images (collate/check.py, witness 'ocr'):
%  1132 disagreements at the first collation.  check.py's own rules settled 96; rules.py settled ~330 more, each rule
%  switched on only after a BLIND PILOT of 107 crops (reader A, 21/22 controls) had measured it on this book: tesseract's
%  glyph faults 8/8, aa read an 8/8, accents it adds 8/8, the layer's glyph faults 7/7 and 8/8, o/ø by the ø rule 8/8, the
%  printer's space before a mark 8/8 (house rule), and the narrow word-division rules; the quote, dash, note-number and
%  word-seam rules failed the pilot and stayed off.  The rest were read at the image BLIND (sheets.py: 579 word crops on
%  58 sheets, two planted controls a sheet, 116 of 116 controls right; six readers) and about 110 spots in the main
%  session from line crops (zoom.py).  STRUCTURE read at the image (reader K): the essays' headings, the volume's title
%  page, every footnote whole (the notes were collated where they stand and then moved to their marks), the signatures.
%  LETTERSPACING (the print's emphasis, \emph): candidate lines by box width (spacedwidth.py; the gap measure is noise on
%  this scan): 163 lines read blind (22/22 roman controls), doubtful ones read in the main session.  ITALICS by stroke
%  slant (check.py --italics): 24 crops read blind; the only italics are Latin quotations.
%  SWEEPS after OPEN 0, each spot read at the image: the W-shortens / T-longer list of the blind decisions (found two
%  W decisions that had dropped a note's mark and a closing «, and one that had joined „der f.“); --moved and --splits;
%  the split sweep; the reader-line sweep; paragraph starts against tesseract's indents; quotation marks (the print
%  closes with « and often with », both as printed); the unknown-word sweep (31 spots: 16 corrections).
%  TRANSLATION as a last reading: the twelve translators listed 79 odd spots (ebook-method/for-idee-1869/tr/odd-*.tsv),
%  ~60 read at the image (zoom zo, zw, zx): mended -- f's whose head did not ink, read t by both OCRs (pp. 21, 36, 309, 319,
%  441; a recurring defective sort, as in Folkelige Foredrag), two blind decisions that had taken a reader's t for the
%  print's f (pp. 322, 539), a speck read as a stop (p. 589), a lost dash (p. 525), the verse on p. 414 set as one quotation,
%  false paragraph breaks at page turns left by the notes' removal (pp. 430, 553, 620, 695, 703); the rest are the print's.
%  SAMPLED PAGE CHECK (pagecheck.py, seed 202610136): pp. 12, 418, 521, 694 and the title page read whole and blind, 1069
%  words: 0 word and 0 punctuation errors; one letterspaced passage in a note missed (p. 694), mended.
%  What this cannot see: a misprint both OCRs normalise into a real word.
%  The ø rule (Hans 2026-10-08): o-forms that are no word (Oieblik, Ifolge, imode ...) are set with ø.
%  DEFECT, letter supplied p. 309: „af“: the f did not ink its head; letter supplied
%  DEFECT, letter supplied p. 319: „fin“: the f did not ink its head; letter supplied
%  PRINTED AS IS p. 418: „Indvendiger“ for Indvendinger
%  PRINTED AS IS p. 423: „aandige“ (so also p. 618)
%  PRINTED AS IS p. 429: „historiske-Virkelighed“, with a hyphen
%  DEFECT, letter supplied p. 441: „finde“: the f did not ink its head; letter supplied
%  DEFECT, letter supplied p. 301: „saaledes“: the second e is broken (a small curved mark); letter supplied
%  DEFECT, letter supplied p. 319: „efter“: the f did not ink its head (both OCRs read t); letter supplied
%  PRINTED AS IS p. 323: „derved. at“: a stop before a lower-case word
%  PRINTED AS IS p. 410: „Læreanstalt. kan“: a stop before a lower-case word
%  PRINTED AS IS p. 410: „Emter“
%  PRINTED AS IS p. 422: „Dömmekraft“, set with ö
%  Not transcribed: running heads, folios, sheet signatures („For Ide og Virkelighed.  20“), the library's stamp; the large
%  initial of each essay is set as a plain letter.  Quotation marks as printed (» « in the text, „ “ in the notes).  The notes
%  are numbered ¹) per page, as printed (p. 553: *)).
%
%  TOOLS (ebook-method/for-idee-1869/, copied from folkelige-foredrag): draft.py, ocr.py, rules.py, sheets.py + reads.py +
%  READER-BRIEF.md, heads.py, spacedwidth.py + spacedlines.py + spaced.tsv, zoom.py, pagecheck.py + PAGECHECK-BRIEF.md,
%  patches.py (PATCHES, NOTES, POST), decisions.tsv, the sweeps, split5.py, rebuild.sh.
% =====================================================================
\documentclass[12pt,a4paper,openany]{book}
\usepackage[utf8]{inputenc}
\usepackage[T1]{fontenc}
\usepackage{libertinus}
\usepackage{libertinust1math}
\usepackage[danish]{babel}
\usepackage[protrusion=true,expansion=true]{microtype}
\usepackage{setspace}
\usepackage[margin=1.3in]{geometry}
\usepackage{fancyhdr}
\usepackage{hyperref}
\hypersetup{pdftitle={To Slags Undervisning}, pdfauthor={Rasmus Nielsen}, hidelinks}
\pagestyle{fancy}
\fancyhf{}
\fancyhead[LE,RO]{\thepage}
\fancyhead[C]{\textit{R. Nielsen: To Slags Undervisning (1869)}}
\setstretch{1.3}
\usepackage{marginnote}
\renewcommand{\thefootnote}{\arabic{footnote})}   % the print numbers each page's note ¹) (p. 553: *))
\newcommand{\opage}[1]{\marginnote{\footnotesize\textit{[#1]}}}
% the essay's heading: title (and subtitle), centred, a short rule; the print's large initial is not reproduced
\newcommand{\essayhead}[2]{\begin{center}{\large #1}\ifx&#2&\else\\[0.3em]{\normalsize #2}\fi\\[0.5em]\rule{3em}{0.4pt}\end{center}\medskip\noindent}
\newcommand{\signature}[1]{\par\smallskip\hfill #1\hspace*{2em}\par}

\begin{document}

% --- front matter: the volume's title page (PDF 7), read at the image (reader K, F7; the page check reads it again blind).
\begin{titlepage}
\begin{center}
\vspace*{3em}
{\Huge For Ide og Virkelighed.}\\[1em]
\rule{8em}{0.4pt}\\[1em]
{\large Et Tidsskrift,}\\[0.8em]
{\small udgivet}\\[0.4em]
{\footnotesize af}\\[0.6em]
{\large R. Nielsen,}\\[0.4em]
{\large Bj. Bjørnson og Rud. Schmidt.}\\[1em]
\rule{8em}{0.4pt}\\[1em]
{\large Aargang 1869.}\\[1em]
\rule{8em}{0.4pt}\\[4em]
Kjøbenhavn.\\[0.4em]
Chr. Steen \& Søns Forlag.\\[1.5em]
{\small Louis Kleins Bogtrykkeri.}
\end{center}
\end{titlepage}

% ===== TEXT (pp. 297-323, pp. 409-441) =====
\essayhead{To Slags Undervisning.}{I.}
% --- p. 297 ---
\setcounter{footnote}{0}%
\opage{297}Undervisning og Dannelse kan anlægges og ledes i
meget forskjellige Retninger; man vilde ved Forsøg paa
at opregne alle mulige Methoder komme til et uoverskueligt
Tal. Her skal kun gjøres opmærksom paa to
i Anlæg, Formaal og Fremgangsmaade uligeartede Skoler,
den \emph{lærde} og den \emph{folkelige.}

Karaktermærket for den lærde Undervisning er Meddelelse
af grundigt sammenhængende Kundskaber med
tilsvarende Aandsdannelse; Hukommelse og Forstand ere
de Sjæleevner, paa hvis Udvikling det fornemmelig kommer
an: her virkes oprindelig paa Menneskets intellektuelle
Væsen, ikke paa det hele Menneske. Opgaven
for den folkelige Undervisning er derimod en personlig
Dannelse, Indvirkning paa Sindet, Vækkelse af Sjælekræfterne
med tilsvarende Oplysning om Menneskelivets
Maal og Vilkaar for den Enkelte, for Samfundet, for
Slægten: her virkes fra Grunden af paa det hele Menneske.
I den lærde Undervisning er og maa Kundskabsinteressen
være det Høieste, medens Folkeoplysningens
oprindelige Drivfjeder er at søge i Livsinteressen. Som
Interesserne ere, maae Bestræbelserne være. Det er
% --- p. 298 ---
\setcounter{footnote}{0}%
\opage{298}nødvendigt at skille de to Slags Undervisning tilbørligt
fra hinanden, naar man vil gjøre sig Rede for Principerne;
men dermed skal dog ikke negtes, at der imellem
begge maa gives mangfoldige Bindeled, at Kundskabsundervisningen
f. Ex. kan udøve en, om end middelbar,
saa dog betydelig Indflydelse paa Karakter og
Livsretning, ligesom at Folkeoplysningen igjen fra sin
Side maa, for at naae nogen Høide, baade forudsætte
og fremme Erhvervelse af fornødne Forkundskaber.

Den lærde Undervisning er hos alle civiliserede Folk
i fuld Gang og fuld Anerkjendelse. Kundskabsskolen
er gammel; den har gjennem Aarhundreder vidst at
skride frem med Tiden, reformere sig selv og trods
mangfoldige baade indvortes og udvortes Hindringer tilkæmpe
sig Udbredelse i stedse videre Kredse. »Kundskab
er Magt«, denne Sandhed staaer klart for alle
Dannede, ligesaa klart for den industridrivende Borger,
der praktisk seer paa Kundskabernes Anvendelse og
Nytte, som for den theoretiserende Forsker, hvis ideale
Maal er Kundskaben selv. En saadan Anerkjendelse er
endnu langtfra bleven den folkelige Undervisning til Del.
Denne Undervisning er da heller ikke Aarhundreder
gammel; den er tvertimod meget ung, endnu i sin første,
famlende Begyndelse; Livstanken i dens indre Eiendommelighed
er --- naar vi nu blot kunde sige det uden
altfor smertelig Forargelse for dem, der have Dispositionen
--- allerførst brudt frem med de Grundtvigske
Høiskoler.

Før vi imidlertid forsøge en nærmere Opklaring af,
hvad og hvormeget der ligger i denne Paastand, maae
vi først afværge nogle Misforstaaelser, som i den nyere
% --- p. 299 ---
\setcounter{footnote}{0}%
\opage{299}Tid, vel særlig i Anledning af den vaagnende folkelige
Stræben, ere fremkomne angaaende det lærde Undervisningsvæsen.

Man har fra et noget almueagtigt Synspunkt været
tilbøielig til at betragte Universitetet og de lærde Skoler
med Alt, hvad derunder hører, som forældede Indretninger,
Institutioner, hvis nær forestaaende Opløsning maatte
ansees som nødvendig Betingelse for friere Fremskridt i
Aanden; man har i Klager over »Bogaristokratiet«, »det
lærde Laugsvæsen«, »Levebrødsvæsenet« o. s. v. gjort
gjældende, at Lektielæsning, Stiløvelser og streng Examen
forvandlede Skolen til en Art Tugthus, en officiel
Tvangsanstalt, der kuede Ungdommens Mod, sløvede den
medfødte Lærelyst, svækkede Livskræfterne, bragte Iskulde
ind i de varme Hjerter og paadrog dem en aandelig
Forkjølelse. Endskjøndt nu mange af disse Indsigelser,
naar vi dømme efter den ydre Virkelighed, ikke just ere
saa ganske ubeføiede, saa ere de dog, naar vi see paa
Ideen, grundløse.

Kundskabsskolen maa til Løsning af sine Opgaver
nødvendigvis gaae ud fra bestemte Forudsætninger, og
stille særegne Fordringer. Den første Betingelse er, at
Eleven indprenter sig, hvad der forklares og foreskrives
ham, ved at lære visse Lektier; anden Vei til Kundskabers
Erhvervelse gives ikke. Nu kan jo vel en saadan
Lektielæsning, især naar Lektierne efterhaanden
blive flere og større, mangen Gang føles som en trykkende
Byrde; her kan da være Tale om at lette den
Unge Byrden, hvis den bliver ham for tung, men her
kan fornuftigvis ikke være Tale om at fritage ham for
Byrden. Byrde tyder paa Arbeide, og den Unge skal
% --- p. 300 ---
\setcounter{footnote}{0}%
\opage{300}lære at arbeide. Som nu Skolen i det Hele vil vaage
over, at Eleven uden Overanstrengelse sættes istand til
at arbeide med Alvor, Dygtighed og Lyst, saa vil den
ogsaa ved omhyggelig Overhøren, Veiledning og Tilrettevisning
vide at kontrollere Arbeidet. Eleven kommer
paa denne Maade til en klar Forestilling om, hvad han
kan, og hvad han ikke kan; han lærer, hvad det er at
lære Noget; han indseer efterhaanden tydeligere, hvor
nødvendigt det er, at kjende Noget tilgavns, at kjende
det i Sammenhæng, grundigt, sikkert og fuldstændigt.
Saa kommer han mere og mere ind under Tankens Tugt
og føler Vægten af gyldige Grunde --- et uundværligt
Element i Karakterens Udvikling; thi det Menneske, der
ikke har Agtelse for klare Tanker og gyldige Grunde,
mangler noget Væsenligt i at være Menneske, hvor levende
Følelse, hvor varmt et Hjerte han end for Resten kan have.

Som en Anke mod det lærde Skolevæsen i dets
nuværende Form høres ikke sjeldent den Indvending, at
Programmet har formange Rubrikker, at der kræves altfor
mange uensartede Ting, saa at Eleverne lære Lidt
af Alt uden at lære Noget tilgavns. Vi skulle her ikke
undersøge, hvorvidt der muligvis for Tiden fordres vel
mange Fag, om disse Fag ikke maaske staae for
udvortes sideordnede, istedenfor at underordnes under
enkelte Hovedfag, der kunde give det Hele større Hold
og Sluttethed; vi skulle kun gjøre opmærksom paa,
at Kundskabsskolen maa sikkre sit intellektuelle Grundlag,
og at dette Grundlag bliver for smalt og Synskredsen
for snever, naar Læregjenstandene ere for
faa og ensformige. Den videnskabelige Forberedelsesskole
kan ikke begynde som Fagskole; just derfor
% --- p. 301 ---
\setcounter{footnote}{0}%
\opage{301}er det, at den optager saa mange uligeartede Fag. De
mange Fag betyde, at det ikke særlig kommer an paa
noget enkelt for sig, at det enkelte Fag kun er et enkelt Element
i den hele Dannelse. Nu er det jo vistnok ganske naturligt,
at blandt flere opvakte og begavede Elever vil En
have Blik og Forkjærlighed for et Fag, en Anden for
et andet, En for Historie, en Anden for Mathematik, en
Tredie for Sprogene o. s. frmdl.; slige Yttringer af individuel
Forkjærlighed kunne endogsaa regnes blandt de
meget lovende Tegn. Men kan end Skolen til en vis Grad
see igjennem Fingre med Enkeltes intellektuelle Tilbøieligheder,
saa tør den dog af Hensyn til Evnernes ligelige,
saavidt muligt, harmoniske Udvikling, ikke give ungdommelige
Liebhaverier utidig Næring, saa lidt som den tør
godkjende den Undskyldning, at et enkelt Hoved skulde
saa udelukkende være organiseret for en enkelt Undervisningsgren, at
det paa Grund heraf umulig kunde sætte sig ind
i Begyndelsesgrundene til de andre. Naar da Fordringen er
af »de mange Fag at lære Noget tilgavns«, saa kan
dette fornuftigvis kun forstaaes saaledes, at Eleven lidt   % PRINTER'S DEFECT (p. 301): „saaledes“: the second e is broken (a small curved mark); letter supplied
efter lidt kommer til Erkjendelse af, hvor vigtigt det
er ham at sammenarbeide det Forskjelligartede, saa at
Overskuelsen over det Hele staaer i det rette Forhold
til Indsigten i dets Enkeltheder; det er ideelle Totaliteter,
et Rige af Love, en Verden af Erkjendelser, som
efterhaanden skal opgaae for hans indre Syn. Det saameget
omtvistede Examensvæsen vil da, naar det sees i
denne Belysning, netop bidrage Sit til Øiemedets Opnaaelse.

Det er en stor og betydningsfuld Opgave, den Studerende
maa løse, naar han grundig skal forberede sig
til en lærd Examen. Det er ikke gjort dermed, at han
% --- p. 302 ---
\setcounter{footnote}{0}%
\opage{302}propper en Mængde Kundskaber og Læresætninger ind i
sin Hukommelse; han maa forstaae at bearbeide det rige,
høist uligeartede Stof, indtil han faaer den hele Masse
organiseret. Hertil udkræves Mere end Forstandsmekanik,
Hukommelse og Flid, hertil udkræves først og sidst
en vaagen, virksom, selvstændig arbeidende Tænkning.
Vel kan Tænkningen endnu ikke vise sin Selvstændighed
derved, at den Studerende paa egen Haand opdager
enten nye Kjendsgjerninger eller nye Sandheder; han er
ligefrem henvist til at modtage og tilegne sig det Overleverede.
Men see vi nærmere til og besinde os paa,
hvad heri ligger, da opdage vi let, at de Kræfter, der
væsenlig betinge videnskabelig Produktivitet, allerede
her ere i fuld Virksomhed. Det er, som bekjendt, ingenlunde
nogen let Sag at sætte sig fuldstændig ind i
en given Lære, selv om Læresætningerne ere fremstillede
i klar systematisk Sammenhæng. Vanskeligheden forøges,
naar flere indbyrdes modstridende Anskuelser skulle
bedømmes og det Stridige udlignes; der hører Aandssmidighed,
Selvbesindelse og Resignation til for at tage
de fremmede Tanker i Besiddelse, leve sig ind i dem og
gjøre dem til sine egne. Øvelsen heri er saa langt fra
at svække den videnskabelige Selvtænkning, at den tvertimod
afgiver en af dens første ufravigelige Betingelser.
Hvormegen umoden Kritik, hvormange løse, frække, kaade
Overhalinger af dybere gaaende Arbeider bringes ikke
til Torvs paa det literære Marked alene af den Grund,
at de Kritiserende have forsømt hin gavnlige Øvelse og
desaarsag mangle den Selvkritik, der skulde advare dem
imod at tage fat paa Opgaver, de saa lidt have gjennemarbeidet,
at de endog, hvad den allertarveligste Rede%
% --- p. 303 ---
\setcounter{footnote}{0}%
\opage{303}gjørelse angaaer, vilde komme ynkelig derfra, dersom
de skulde bestaae en offenlig, upartisk, men streng
Examen.

Det Næste, Forberedelsen kræver, er at faae de
uligeartede Erkjendelser grupperede i Bevidstheden saaledes,
at Detaillerne træde tydelig frem uden at skade
Overblikket. Under denne Bestræbelse er Enhver i
strengeste Forstand henvist til sig selv; thi den Fremgangsmaade,
som passer for den Ene, kan være høist
ubekvem for den Anden. Den psykologiske Vexelvirkning
mellem Fantasi, Forstand og Hukommelse, hvorpaa
det her kommer an, kan i de indre, fine, ofte umærkelige
Overgange være høist forskjellig hos de Forskjellige.
Den, der f. Ex. har en meget omfattende og tro Hukommelse,
vil være tilbøielig til at begynde med Enkelthederne og
betragte det Hele som en Sum af mangfoldige Dele; først
ved at huske hjælpes han til at tænke: Forstanden opererer
paa Hukommelsens Grundlag. Den, hos hvem Forstand
og Tænkning har en naturlig Overvægt, maa gaae den
omvendte Vei. For ham er det allerførst af Vigtighed at
fremhæve det Almene, trænge ind i Sætningernes Sammenhæng
og bevare Overskuelsen. For ikke at bebyrdes
med overflødigt Stof, maa han ved hvert Skridt
fremad lade Tænkningen strengt prøve og nøiagtig afveie,
hvilke og hvormange Enkeltheder der i den givne
Sammenhæng nødvendig høre med til Forstaaelsen af
det Hele: det er Forstanden, som ved denne Økonomi
kommer Hukommelsen til Hjælp. At nu begge, baade
han med den stærke Hukommelse og han med den indtrængende
Forstand, maa stræbe til eet og samme Maal,
er en Selvfølge; men det gjør, hvad Arbeidsmaaden an%
% --- p. 304 ---
\setcounter{footnote}{0}%
\opage{304}gaaer, en væsenlig Forskjel, om det er Hukommelsen,
der oprindelig bærer Forstanden, eller omvendt.

Videnskabelig Gjennemgang gjennem en stor, i sig
selv afrundet Lærebygning, hvis Portaler og Søiler, hvis arkitektoniske
Led og Dele staae i symmetrisk Samklang med
hverandre indbyrdes og med det Hele, og derhos et tilsvarende
psykologisk Indblik i Selvets intellektuelle Væsen,
et grundigt Bekjendtskab med Ens egen aandig-individuelle
Natur, en Selvkritik, der afslører Mangler og
Fortrin og afgiver Maalestok for fremtidig Arbeidsdygtighed,
ere Fordele, som neppe Nogen, der ved en grundig
Examens-Forberedelse har holdt Anstrengelsen ud,
vil fortryde at have erhvervet sig.

Der kan da, saasandt vi agte Videnskaben og fastholde
det videnskabelige Synspunkt, ikke gjøres nogensomhelst
begrundet Indvending imod Skolen som Kundskabsskole;
det er meget mere at ønske og haabe, at
den lærde Undervisning, uanfægtet af Tidens vexlende
Stemninger, maa i Eet og Alt blive sin Karakter tro.
Derimod er der i de herskende Forestillinger om Kundskab
og Dannelse en Grundvildfarelse, som paa det
Kraftigste maa modarbeides.

Det er en hos Mange dybt indgroet Fordom, at Kundskabernes
Mangfoldighed er Dannelsens Kjendemærke,
at den, der har mange Kundskaber, kan tale mange
Sprog og tale med om Alt, er den sande Dannede, medens
den, der har færre Kundskaber og kun taler eet
Sprog, røber Mangel paa Dannelse og kan som indskrænket
Individ let blive overseet i det dannede Selskab.
At denne Synsmaade, der staaer i nøie Forbindelse
med Tidsalderens herskende Tilbøielighed til at
% --- p. 305 ---
\setcounter{footnote}{0}%
\opage{305}sætte Menneskets udelukkende Fortrin i det Intellektuelle,
er til ubodelig Skade for Aandslivet i alle Retninger,
vil ved lidt nærmere Overveielse let kunne indsees.

Den er for det Første til Skade for den lærde Undervisning;
thi den bevirker, at Kundskabsskolen idelig
staaer Fare for at falde ud af sin oprindelige Karakter.
Er det Lærdom og Videnskab, der væsenlig skal betinge
Aandsudviklingen, da maa Skolen paa det Strengeste
indskærpe, at disse høie intellektuelle Goder allerførst
eftertragtes for deres egen Skyld uden nogetsomhelst
Hensyn til udvortes, endelige Fordele. At lade
sig optage i den lærde Skole er da det Samme som at
indvie sit Liv, sin bedste Tid og sine bedste Kræfter til
Videnskabens Dyrkelse. Vistnok bør der fra Statens
Side vaages over, at Videnskaben altid maa kunne sikkre
sine sande Dyrkere en virkelig Livsstilling, en hæderlig
Plads i Samfundet; men slige lovende Udsigter tør ikke
komme i Betragtning, naar Spørgsmaalet er om de første
Bevæggrunde. Før Videnskaben kan tjene som Middel,
maa den først være Maal i sig selv, tillægges Værd
i og for sig selv og elskes for sin egen Skyld. Den
Discipel, som, idet han træder ind i Skolen, ikke medbringer
en ungdommelig stor, om end ungdommelig uklar,
Forestilling om det Herlige, det Ophøiede i at stræbe
efter Kundskab og Indsigt, ham ville Muserne ikke hilse:
han træder ind som en Uværdig; og den Discipel, som,
uagtet Skolen anvender sine bedste Lærekræfter paa at
aabne hans Øie og vække hans Sands, dog ikke kan
drive det videre end til det praktiske Spørgsmaal, hvad
Nytte man dog har af at lære saamange unyttige Ting,
bør, om han end for Resten læser Lektierne nok saa
% --- p. 306 ---
\setcounter{footnote}{0}%
\opage{306}upaaklageligt, meddeles consilium abeundi, thi han er
ikke paa sin rette Plads, han er gaaet feil. Men dersom
Skolen paa denne Maade skulde opretholde de ideale
Fordringer, da følger det af sig selv, at den, for at
vinde i Indhold, maatte tabe betydeligt i Omfang. For
at blive ved den klassiske Dannelse, vilde to eller tre
Latinskoler, som hver aarlig i Gjennemsnit sendte ti a tolv
Studerende til Universitetet, vel omtrent være Maximum
i et saa lille Land som Danmark. Nu derimod, da uklare
Forestillinger om Aandslivets Hindringer og Betingelser
for en Del gjøre Udslaget, staaer Sagen anderledes.

I den skikkelige Mening, at den studerende Vei dog
altid er en sikker og hæderlig Vei, at en Skat af Kundskaber
er en kostelig Skat, og at »det Bedste, man kan
gjøre for sine Børn, er at give dem en god Opdragelse«,
see Forældre og Værger sig bekymrede om efter høiere
Dannelsesanstalter, og see da ofte feil; thi det gaaer
en masse løs paa Latinskolerne, Gymnasierne og Universitetet.
Dette tyder, især naar vi see hen til en
kommende Tid, hvori neppe nogen velstaaende Borger
vil have Aarsag til at misunde Embedsstanden det
gyldne Overflødighedshorn, aabenbart paa en dybere liggende
Grundvildfarelse. Det kan umulig være Hensyn
til Fordelen eller de gode Udsigter, som her har Overvægten;
det maa være noget Andet. Det er den taagede
Blanding af to Slags Dannelse, den intellektuelle og den
personlige, der forvirrer Opfattelsen. --- Realskolen f. Ex.
er utvetydig en Kundskabsskole og intet Andet. Den,
der vil sætte sin Søn i en Realskole, bliver sig let bevidst,
hvad det er, han egenlig vil. Han vil, at det
unge Menneske skal arbeide sig frem, skal erhverve sig
% --- p. 307 ---
\setcounter{footnote}{0}%
\opage{307}en Sum af nyttige Kundskaber i bestemte Fag. Han
kan for Resten have den største Agtelse for Videnskabelighed,
særlig for Realvidenskaberne; han kan forstaae,
at en Arbeiden med Hovedet styrker Evnerne og giver
den Duelige en egen Forstandsoverlegenhed; men han
bedaares ikke let af uklare Fantasier angaaende det
Ideale i Undervisningen selv: at Realskolen just i særegen
Grad skulde forædle Hjertet og forhøie det Aandelige
i Personlighedens Udvikling, falder ham neppe
ind. Anderledes med den klassiske Undervisning; den
synes nu engang at have faaet Privilegium paa at fremelske
det Ædleste og Dybeste i den menneskelige Aand,
saa at den baade meddeler nyttige Kundskaber og tillige
giver de unge Sjæle den finere aandelige Næring,
der er Betingelsen for ideal Personlighed. Men netop
dette delte, dobbeltsidige Hensyn, Hensynet dels til de
nyttige Kundskaber, dels til den finere aandelige Næring,
medfører en Tvetydighed, der er yderst farlig for
Skolen. Den tør nu ikke længere betragte sine Elever
udelukkende som en Kreds af unge vordende Videnskabsdyrkere;
langtfra! her er jo Mange tilstede, som vitterlig
aldrig have tænkt paa at offre deres Liv til Videnskaberne;
de ere her blot for Kundskabsindsamlingens Skyld
og for at faae nogen klassisk Dannelse. Endnu mindre
tør Skolen lægge det afgjørende Eftertryk paa Kundskabernes
Nytte; thi hvad det Nyttige angaaer, kan den
umulig konkurrere med Realskolen. Skal Latin og Græsk
læres for Nyttens Skyld, skulle Herodot og Homer, Cicero,
Virgil, Ovid og Horats læses i Grundsprogene
udelukkende for Kundskabernes Skyld og for den praktiske
Anvendelses Skyld, da seer det sørgeligt ud med den finere
% --- p. 308 ---
\setcounter{footnote}{0}%
\opage{308}Næring, den klassiske Dannelse og den forventede Sjæleadel.
Saa staaer Skolen i al Tvetydighed som en
høiere Dannelsesanstalt --- med mange Karakterbøger
--- uden Karakter. Med denne Tvetydighed kan det
nemlig ikke undgaaes, at de tilstrømmende Elever jo
komme til at udgjøre en blandet Befolkning. Omkring
de virkelig Kaldede, Minervas ægte Sønner, samler sig
efterhaanden en større Flok af uægte Sønner, af Saadanne,
som, da de nu engang skulle gaae »den studerende
Vei«, vistnok paa Afstand have en vis Respekt for den
høie Gudinde, og vel ogsaa med Tiden ville gjøre hende
den Ære at anbringe hendes Ugle paa Kaskjetten; men
som for Resten, hvad det Videnskabelige angaaer, rent ud
sagt, ikke vide med sig selv, hverken hvad det Hele
skal betyde, eller hvad de i aandelig Forstand egenlig
have at gjøre. Stemningen hos Flertallet vil da udøve
en Indflydelse paa Aanden i Skolen, der ikke er til det
Bedre. En dygtig Bestyrelse i Forening med nidkjære,
duelige Lærere kan vel afgive Sikkerhed for, at Sædelighed,
Orden og taalelig Flid i det Hele opretholdes;
men den Opvakthed, der betinger dybere Forstaaelse,
den ideale Sands, den aandige Hunger og Tørst lader
sig ikke indblæse. Som det er modbydeligt at skulle
spise og spise og smage paa mange Retter, naar man ingen
Appetit har, saaledes maa det ogsaa falde den Unge
besværligt at skulle læse og læse og lære de hundrede
forskjellige Ting, af hvilke maaske neppe mere end een
eller to i Virkeligheden og naturligt kunde, --- »hvis
han blot havde Tid« --- interessere ham.

Jo mere de indvortes Hensyn aftage, desmere voxe
de udvortes. Det er ikke Aanden i de Gamle, ikke de
% --- p. 309 ---
\setcounter{footnote}{0}%
\opage{309}store Verdenstanker, ikke Sprogets Rytmer, ikke den
klare klassiske Form, der fængsler Eleven og tilskynder
ham til at lære sit Pensum; nei, det er andre mere udvortes
og reelle Hensyn, Hensyn til »mg og ug« og til
»Oprykning ved næste Examen«. Aandsarbeidet faaer paa
denne Maade Lighed med et Fabrikarbeide; man besørger
det Nødvendigste af Hverdagsgjerningen mekanisk
og med døsig Flid, men arbeider paa Kraft i den travle
Tid; man trykkes bestandig kun af een Bekymring, en
stille Bekymring, der voxer, alt som Afgjørelsens Stund
nærmer sig, Bekymringen for at »bestaae ved Examen«.
Under Trykket af Skoletvangen og den i mange Maader
besværlige Lektielæsning er der imidlertid et Noget, der
bevæger sig i den ungdommelige Fantasi og lyser med
Idealitetens Glands, det er Forestillingen om det frie
Liv ved Universitetet. »I Skolen læse vi kun for Examens
Skyld; først naar vi komme til Universitetet, begynde
vi at studere«. Denne Forestilling vil skuffe
Mange; den vil i Reglen skuffe alle dem, der ikke
alt fra Begyndelsen af have følt, hvad det er at
sætte de indvortes Hensyn høiere end de udvortes.
Thi den første »praktiske« Fordring, Universitetet stiller
til de unge Studerende, er, at de snarest muligt
forberede sig ikke blot til een, men til tre Examina, af   % PRINTER'S DEFECT (p. 309): „af“: the f did not ink its head; letter supplied
hvilke de to første allerhelst maatte afsluttes i Løbet af
det første Aar, for at der saa med udelte Kræfter kan
arbeides paa den store Embedsexamen; »først naar Embedsexamen
er vel overstaaet, kan man egentlig begynde
at studere«. Er allerede i Skoletiden Hensynet til Examen
blevet et udvortes Hensyn, og Forestillingen om, hvad
det er at læse til en Examen, slaae sig igjennem og
% --- p. 310 ---
\setcounter{footnote}{0}%
\opage{310}bestaae ved en Examen, behørigt trivialiseret: hvad Under
da, om den gamle nøgterne, triviale Forestilling
dukker op paany, saa at Fabrikvæsenet med store Lektier,
mekaniske Manuduktioner, Hastværksrepetitioner og
ilfærdige Overhalinger, inden man ret veed deraf, er i
fuld Gang.

Det er Kundskabsskolens medfødte Ensidighed, at
den, naar den samvittighedsfuldt bevarer sit oprindelige
Præg og virker Alt, hvad den formaaer, i Aandens Tjeneste,
dog ikke er istand til at udvikle hele den menneskelige
Aand; faaer det nu desuagtet Udseende af, at
den lærde Høiskole er den eneste aandelige Dannelsesanstalt,
der kan gives, saa maa den synke fra sit oprindelige
Niveau ned til at tjene det blotte Udvortes, til
Asyl for det tomme Skin, et dødt mekanisk Lærdomsvæsen, et af endelige Interesser drevet Fabrikvæsen, for
hvis Kvalme Muserne flygte.

Er den uklare Forestilling om »Dannelse«, »Aandsdannelse«,
»høiere Dannelse«, allerede til Skade for det
videnskabelige Undervisningsvæsen, da er Skaden endnu
langt større, naar vi stige nedad og betragte Undervisningen
i de Kredse af Livet, hvor der ikke kan være
Tale om Videnskab. Alle vore offenlige Skoler, Borgerskolerne
og Almueskolerne, ere jo uden Undtagelse
Kundskabsskoler, alle anlagte paa Lektielæsning, Udenadslæsning
og Examenslæsning, allesammen indrettede
efter det lærde Forbillede. Dersom nu disse Skoler alle
bleve opfattede under Synspunktet af populære Realskoler,
saa kunde herimod ikke være Noget at indvende;
thi nyttige Kundskaber ere nødvendige for Alle.

Men saaledes forholder det sig ikke. I Borger- og
% --- p. 311 ---
\setcounter{footnote}{0}%
\opage{311}Almueskoler lægges der udtrykkelig an paa at vække
det høiere Liv i den unge Slægt, at virke opdragende.
Til Opdragelsen hører da ikke blot, at Børnene lære at
læse, skrive, regne o.s.v., dertil hører fornemmelig, at
de undervises om det for Livet Høieste og Vigtigste, at
de undervises i Religion. Betragte vi nu Religionsundervisningen
i disse Skoler lidt nærmere, da møde vi
den ovenfor berørte Tvetydighed igjen, men under en
endnu bedrøveligere Skikkelse. At meddele Religionskundskaber
og at vække det religiøse Liv, disse to uligeartede
Opgaver fremkalde to Bestræbelser, af hvilke den
ene uophørlig modarbeider den anden. »Timen i Religion
er en Slags opbyggelig Time; de religiøse Sandheder
kunne ikke meddeles enten som blotte Forstandslærdomme
eller som udvortes historiske Kjendsgjerninger. Derfor
kunne jødiske Børn ikke undervises sammen med kristne;
de kunne have alle de andre Timer fælles, kun ikke Timen
i Religion. Religionens Lærdomme maae komme
fra Hjertet og gaae til Hjertet, de maae derfor meddeles
af en troende Lærer, som paa bedste Maade søger
at vække og befæste Troen hos de Unge.« Dette er Bestræbelsen
i den ene Retning. Men hermed forbindes
nu en anden Bestræbelse, og den er for Skolen som
Skole endnu vigtigere, det er Bestræbelsen for at meddele
Kundskaberne i et vist befalet Omfang.

Hvad er Religion? »Man kalder det Religion at
have Kundskab om Gud og hans Villie, samt om den
Maade, paa hvilken han vil æres og dyrkes af os«; disse
Begyndelsesord i den ældre befalede Lærebog ere i høi
Grad betegnende for det Synspunkt, hvorfra Religionsundervisningen
opfattes i Skolen. Ogsaa i Religionen
% --- p. 312 ---
\setcounter{footnote}{0}%
\opage{312}er der Lektier, ogsaa i Religionstimen skal der høres
paa Lektien, ogsaa i Religionen skal der tages Examen.
Hvad er det nu, der kan prøves ved en Examen? Ikke
Troeslivet selv, eller hvorvidt de Unge hver især virkelig
have Troen. En saadan Examen i Troen har Tilværelsens
Herre udelukkende forbeholdt sig selv at anstille.
Det Eneste, der kan prøves, er Kundskaberne.
Men det, der ved Examen udelukkende er Prøvelsens
Gjenstand, er ogsaa Skoleundervisningens høieste, vigtigste
Maal, og Bestræbelsen for Kundskaberne faaer da
af denne Grund en naturlig og nødvendig Overvægt over
den vækkende Bestræbelse, Bestræbelsen for Livet.

Grundlaget for den offenlige, officielle Religionsundervisning
er en justeret Lærebog, et kort Udtog af den
kirkelig systematiske Dogmatik: her have vi det lærde
Forbillede. Lærebogen med dens Skriftsteder er bestemt
til Udenadslæsning, især gjælder det om at kunne Skriftstederne
paa sine Fingre. Udenadslæsning er for Børn
en aandelig mekanisk Syssel; saa begynder det med
Aandsmekanik; men Aandsmekanik i Religion slaaer
Livet ihjel. Det nytter ikke, at Lærebogen udtrykkelig
siger, at Religionen ikke alene er en Kundskab, men
ogsaa et Liv; thi, naar det staaer i Lektien, gaaer det
øieblikkelig over i Udenadsmekaniken, at Religionen er
et Liv; det nytter ikke, at der autoriseres en ny Lærebog,
som begynder lidt mere hjertelig med det Spørgsmaal:
»Hvad skal fremfor Alt ligge os paa Hjerte?« thi,
hvad der »fremfor Alt skal ligge os paa Hjerte«, løber
uden videre med i Ramsen. Dersom der blev udsat en
Prisbelønning for Besvarelsen af det Spørgsmaal: »Hvorledes
skal man under Skin af at ville indprente Børn
% --- p. 313 ---
\setcounter{footnote}{0}%
\opage{313}den aabenbarede Religions Hovedlærdomme uden Forargelse
bidrage Sit til, at de fra Begyndelsen af kunne
faae hemmelig Lede til al Religion?« da vilde Opgaven
uden Tvivl bedst kunne løses ved en grundig psykologisk
Udvikling af, hvad der under Udenadslæsning af en
dogmatisk Lærebog halv bevidst, halv ubevidst maa
foregaae i Barnets Sjæl.

Vistnok kan man med nogen Grund beraabe sig
paa, at det ikke er en slavisk Udenadslæren, hvorom
det gjælder, at Lærebogen kun er en Ledetraad til forstandig
Afbenyttelse, og at Læreren altid vil sørge for
at vække Børnenes Eftertanke, idet han forklarer dem
det Læste saaledes, at det baade virker paa Følelsen
og oplyser Forstanden. Vi tvivle ikke om, at den
retsindige Lærer jo vil i begge Henseender gjøre sit
Bedste; vi mene kun, at det i sig selv Modsigende
ikke lader sig udføre. Hvad Følelsen angaaer, er det
umuligt, at Sandheder, der meddeles i Form af rene
Læresætninger og have saa Lidet med Barnets virkelige
Liv at gjøre, just skulde efterlade noget særdeles
Indtryk eller fæste sig dybt i det barnlige Sind. Man
tage f. Ex. en Læresætning som denne: »Ægteskabet
er den inderlige Forbindelse mellem Mand og Hustru,
som Gud har indstiftet til gjensidig Hjælp og Menneskeslægtens
Forplantelse. Ægtefolk skulle leve sammen
i Kjærlighed og Troskab indtil Døden.« Saa sand denne
Sætning nu end i og for sig er, saa vilde det jo dog
være unaturligt, om den virkede stærkt paa den lille
tiaarige Pige eller vakte alvorlige Tanker hos den elleveaarige
Dreng. Hvis den lille Pige og den lille Dreng,
naar de udenfor Skoletiden talte med hinanden om det,
% --- p. 314 ---
\setcounter{footnote}{0}%
\opage{314}der naturlig maatte ligge dem begge paa Sinde, følte
sig uvilkaarlig tilskyndede til at gjennemgaae og forklare
det nævnte Stykke samt derunder hørende Skriftsteder
med hinanden paa fri Haand: saa vilde man vel
finde det noget betænkeligt, ja uhyggeligt, fordi det var
unaturligt. Derimod vilde man finde det ganske i sin
Orden, om de flittige Børn, den lille Dreng og den lille
Pige, i Læsetimer udenfor Skolen overhørte hinanden i
den foreskrevne Lektie og passede nøie paa Skriftstederne,
baade det med »Ægteskabet være hæderligt« og
det med »Mændene ere skyldige« og det med »I Kvinder!
værer Eders egne Mænd underdanige«; det Barnlige
er, at der ikke tænkes videre ved det, at det gaaer
paa Ramse; jo raskere Ramsen gaaer, desto uskyldigere
er det Hele.

Lærebogen er objektiv; dens dogmatiske Sætninger
ere alle indrettede efter Systemet, og det objektive System
kan ikke tage Hensyn til den barnlige Alders endnu
umodne Subjektivitet. Man indvende ikke, at det kun
er enkelte Læresætninger, hvis Anvendelse og rette Forstaaelse
ligger udenfor Børnenes Synskreds; thi det lader
sig let vise, at der ikke er en eneste religiøs Sætning,
som, vel at mærke, i Form af Lærebogssætning er, naar
vi tage det strengt, tilgængelig for den barnlige Forstand.
»Gud er retfærdig, det er: han lader det Gode
faae Fremgang og Løn, og han hæmmer og straffer det
Onde«; denne Sætning kan et Barn ikke forstaae. Tages
den enfoldigt efter Ordene, saa kommer den i Strid
med Kjendsgjerninger, paa hvilke den Unge allerede tidlig
vil blive opmærksom; tages den enfoldigt efter Aanden,
saa bliver Meningen den, at de i Sandhed Gode,
% --- p. 315 ---
\setcounter{footnote}{0}%
\opage{315}hvad jo kan sees af Profeternes, af Apostlenes og fremfor
Alt af Kristi Exempel, netop maae lide i denne Verden,
medens det Onde paa mange Maader faaer »Fremgang
og Løn«; men op til det Aandens Høidepunkt,
hvorfra et saa skærende Misforhold viser sig foreneligt
med Guds Retfærdighed, kan Barnet ikke løftes.
Skal den opbyggelige Fortolkning gaae ud paa at
vise, at det Onde jo vistnok til en Tid kan faae Fremgang
i Verden, men kun til en Tid, saa at man for
Udfaldets Skyld dog i Reglen staaer sig bedst ved at
gjøre det Gode: da bliver Forklaringen ræsonnerende;
der fremkommer et Blandingsbegreb, et Miskmask af
Fromhed og Egennytte, Gudsfrygt og verdslig Klogskab,
som forvirrer Barnets Tanker og nedslaaer dets ideale
Forestilling om Guds Retfærdighed.

»Men« --- siger man maaske --- »denne hele Betragtningsmaade
gjør kun Virkning ved at overdrive
Skyggesiderne og stirre paa den Ufuldkommenhed,
som klæber ved alt Menneskeligt; den overseer Livets
Kjendsgjerninger og staaer i Strid med den sikkre Erfaring.
Thi Erfaring lærer, at skikkelige, sædelige, flittige,
nogenlunde begavede Børn uden Besvær baade lære
Lærebogen og forstaae Forklaringen og finde Opbyggelse
derved. Den, der tvivler herom, behøver blot, naar Leilighed
gives, at træde ind i Skolen og overvære en
Katekisation eller, hvad der endnu er tilforladeligere, være
tilstede ved den egenlige Modenhedsprøve, den sidste
offenlige Examen. Ved at høre paa Examen vil enhver
Billigttænkende kunne overbevise sig om, at Børnene
virkelig have forstaaet Lærebogen, at der ved sammenhængende
Kundskaber i den hele Religion er lagt en
% --- p. 316 ---
\setcounter{footnote}{0}%
\opage{316}sikker Grund, saa at Troen er vis og Forstanden oplyst.
Det er dog en høist vigtig Sag, at de Unges Forstand
paa Religionen tidlig udvikles, saa at de kunne gjøre
Regnskab for, baade hvad de skulle troe, hvormeget de
skulle troe, og hvorfor de skulle troe det Altsammen«.
Vi forstaae meget godt, hvad heri ligger. Det er med
en ofte forbausende Sikkerhed, at opvakte og vel underviste
Børn kunne ved en Skoleexamen svare paa Spørgsmaal,
der ellers sætte graanende Tænkere i Forlegenhed.
»Hvad vil det sige, at Gud har skabt Himmelen og Jorden?«
»Gives der andre fornuftige Skabninger end Menneskene?«
»Ere alle Englene gode?« »Hvad er Guds
Øiemed med Verdens Opholdelse?« paa slige Spørgsmaal
kan det halvvoxne Barn ofte svare saa tydelig, rent og
bestemt, at det beskæmmer de Ældre. Dersom der nu
virkelig var Sandhed i al denne Forstaaelse, saa var
det Vidunderlige naaet. Men Sandheden er, at det Hele
er Skin, et psykologisk, sjælefordærvende Bedrag.

Enhver, som har givet sig noget af med at undervise,
vil være bleven opmærksom paa den Smidighed,
hvormed en ungdommelig Forstand formaaer at efterligne
Lærerens Tanker uden egenlig at tænke dem. Der er
en Tankernes Efterligning, der kan være saa skuffende,
at det livagtig seer ud, som havde Lærlingen virkelig
forstaaet Sagen, idet han endog for en Del med egne
Ord nøiagtig eftergjør en foreskreven Tankebevægelse;
og dog er det en Skuffelse, en som oftest naiv Selvskuffelse. I exakte Videnskaber, Mathematik, Mekanik,
er Skuffelsen endda temmelig let at opdage; men ikke
saaledes i Aandslæren. Man vil i Filosofien f. Ex. kunne
fremstille ethvert Problem, selv det allervanskeligste,
% --- p. 317 ---
\setcounter{footnote}{0}%
\opage{317}paa en saadan Maade, at Begynderen bliver i Stand til
at følge den nøiagtig betegnede Tankevei gjennem alle
Krumninger med en saadan Klarhed, Sindighed og tilsyneladende
Naturlighed, at det tager sig ud, som havde
det unge Menneske tænkt over Sagen i mange Aar; og
dog er det Hele maaske kun en skuffende Efterligning.
Derfor maa der ved den filosofiske Undervisning fremfor
Alt sørges for, at Skemaet ikke kommer til at gjælde
for Systemet, at ikke en stereotyp Lærebog med afmaalte
Kapitler og staaende Paragrafer Aar efter Aar
sætter sig saa fast i Traditionen, at Filosofien selv bliver
en Tankelektie. At der hører en vis Tænksomhed til
for at lære en Tankelektie, er jo vist nok; men det er
ogsaa vist, at der er Intet, som mere svækker Evnen
til Selvtænkning end denne hule, overfladiske Efterligning,
denne tankeløse Tænkning, der tilsidst maa ende i intellektuel
Løgnagtighed.

Er den tankeløse Tænkning en farlig Snyltegjæst
ved den videnskabelige Undervisning, hvor der dog kan
anbringes saamange og forskjellige Korrektiver, da er
den, som Enhver vil forstaae, endnu langt farligere ved
Børneundervisningen, allerfarligst ved Undervisning i
Religion. De objektive religiøse Læresætninger ligge i
den Grad udenfor Børnenes virkelige Fatteevne, at disse
hverken kunne føle eller tænke noget i Sandhed Væsenligt
ved nogen af dem. Man tage f. Ex. Nadverlæren.
»Brødet og Vinen i den hellige Nadvere ere«
--- siger Lærebogen --- »ikke blotte Tegn paa Jesu Legeme
og Blod; heller ikke forvandles Brødet og Vinen
til Jesu Legeme og Blod, men i, \emph{med} og \emph{under} Brødet
og Vinen modtage vi i den hellige Nadvere Jesu
% --- p. 318 ---
\setcounter{footnote}{0}%
\opage{318}Legeme og Blod, som paa en for os ubegribelig Maade
er virkeligen der tilstede«. Vi spørge: hvad Indtryk
skal denne dogmatiske Oplysning gjøre paa den unge
Piges barnlige Sind? Hvorledes skal den halvvoxne
Dreng være i Stand til enten at tænke eller forestille sig
noget Bestemt ved den givne Forklaring? Det Naturligste
er, at han ikke engang forsøger derpaa.
Skulde det et Øieblik falde ham ind at ville forstaae
Forklaringen, da er han ilde faren. »Hvorfra vide vi
Noget om den Maade, paa hvilken Kristi Legeme og
Blod er nærværende i Nadveren? Det maa vel være en
guddommelig Autoritet, der siger os det; men hvilken?
Bibelen kan det ikke være; thi af Indstiftelsens Ord:
\emph{det er mit Legeme; dette er det Nye Testamentes
Kalk i mit Blod,} kunne vi ikke see, hvorledes
det forholder sig med \emph{Maaden.} Der staaer jo
ikke: i, \emph{med og under dette} (Brødet) \emph{er mit Legeme;}
der staaer udtrykkelig, at \emph{dette} (Brødet) \emph{er}
\emph{mit Legeme;} endnu mindre siges der, at Blodet er i
Vinen, thi det ligger nok saa nær at forstaae Indstiftelsesordene
saaledes, at det er \emph{Vinen} (det Nye Testamentes
Kalk), \emph{der er i Blodet.} Have vi nu ikke en
guddommelig Autoritet, saa maae vi vel have en menneskelig.
Det maa være store Religionslærere, lærde og
fromme Mænd, som ved alvorlig Tænkning have udgransket
Maaden, og det saa bestemt og nøiagtigt, at vi nu
lige saa trygt kunne støtte os til deres menneskelige
som til Skriftens guddommelige Autoritet; men er det
muligt? Her læres jo udtrykkelig, at Maaden er ubegribelig;
hvordan skal man da forestille sig, at der er
Mennesker, som nøiagtig kunne fastsætte og bestemme,
% --- p. 319 ---
\setcounter{footnote}{0}%
\opage{319}altsaa begribe en ubegribelig Maade. Her skal baade
tænkes og dog ikke tænkes, baade forestilles og dog
ikke forestilles noget Bestemt ved \emph{den ubegribelige}
\emph{Maade:} hvor besynderligt!«

Det er en stor Lykke for Religionslærerne i Skolen,
at det end ikke falder de bedst begavede Drenge
ind at ville tænke over, hvad der staaer i Lærebogen.
Hvis denne Uskik kom i Gang, vilde det være til megen
Forstyrrelse i Læretimerne; thi selv om Børnene vare
saa sædelig oprigtige, at de aldrig tillode sig at spørge
enten af Selvvigtighed eller efter løse Indfald: saa vilde   % PRINTER'S DEFECT (p. 319): „efter“: the f did not ink its head (both OCRs read t); letter supplied
der dog fremkomme saamange Spørgsmaal, at der --- for
Overhøringens Skyld --- ikke kunde blive Tid til at besvare
Halvdelen, og, hvad værre er, Spørgsmaal, som
vilde sætte Læreren selv i den pinligste Forlegenhed.
Nu derimod, da hin Uskik lykkeligvis ikke finder Sted,
gaaer Sagen sin jevne, rolige Gang. De Unge lære af
Instinkt, thi Instinktet er en fin og skarpsindig Læremester,   % PRINTER'S DEFECT (p. 319): „fin“: the f did not ink its head; letter supplied
at der i de indprentede Forestillinger er Noget, man ikke
skal forestille sig, i Forstaaelsen Noget, man ikke skal forstaae:
saa lære de Unge, ikke at tænke religiøst, nei, at efterligne
og gjengive foreskrevne religiøse Tankelektier, saa kan
Efterligningstalentet med dets glimrende Hulhed komme til
Ære og Anseelse ved Examen. Det er ikke sand Dannelse,
end ikke Begyndelsen dertil; det er en mekanisk
Dressur. Som Rekrutten dresseres til at lære sine Reglementer,
saaledes dresseres den opvoxende Skoledreng
til at lære sine Dogmer; og dog ikke engang saaledes,
ikke nær saa praktisk, thi den unge Soldat vil strax
have Brug for Reglementerne, men Drengen er endnu
langt fra at have Brug for Dogmerne.

% --- p. 320 ---
\setcounter{footnote}{0}%
\opage{320}Blev Livsinteressen virkelig sat over Kundskabsinteressen,
da var det ikke Lærebogen, men Bibelhistorien,
der kom til at afgive Grundlaget. Nu læres der
jo vistnok Bibelhistorie baade i Almue- og Borgerskoler;
men ogsaa i denne Henseende hersker det lærde Forbillede.
Istedenfor at lade de store Aandsbegivenheder
træde frem for Anskuelsen gjennem Bibelens egen barnlig
ophøiede Fremstilling, hvori Modsætningen af Guddommeligt
og Menneskeligt, af det Himmelske og det
Jordiske, den ideale og den triviale Virkelighed kan
føles og fattes uden ræsonnerende Omsvøb, anvendes der
nu til skematisk Oversigt et lille aandsfattigt Udtog af
den hele Religionshistorie, at Barnet ligelig kan faae
Noget at vide om alle de Ting, Gud har gjort fra Verdens
Begyndelse, fra Skabelsens første Dag indtil Rigsdagen
i Kjøbenhavn 1536. Her kan ikke dvæles ved
enkelte store Optrin, ved saadanne tanke- og farverige
Livsbilleder, der kunde give de unge Sjæle et dybt uforglemmeligt
Indtryk af Idealets Herlighed, af Aandens
Magt, naar Gudsstemmen høres, naar det Vidunderlige
skeer, naar det Overnaturlige omstraaler det Naturlige.
For Udtogets Skyld blive alle Begivenheder, de overnaturlige
med de naturlige satte paa samme Virkeligheds
Trin og stillede i samme Lys: at »Herren sendte sin
Engel til Daniels tre Venner i den gloende Ovn«, er
virkeligt, at »Jonas blev slugt af en stor Fisk, der
kastede ham op igjen, saa at han kom levende paa
Strandbredden«, er virkeligt, virkeligt i samme Forstand
som det er virkeligt, at Rhoboam var samtidig med
Jeroboam, at Titus var Biskop paa Kreta og Kristen
Pedersen født i Svendborg. Saa har man et ligeligt
% --- p. 321 ---
\setcounter{footnote}{0}%
\opage{321}Overblik over Alt, faaer i en vis Sammenhæng Lidt at
vide om al Slags Virkelighed, Lidt om Titus, Lidt
om Daniel, Lidt om den gloende Ovn, Lidt om den
store Fisk, Lidt om Jonas og Lidt om Kristen Pedersen.

Det nytter naturligvis ikke at ville ræsonnere med
Børn om Forskjellen mellem det Sande og det Virkelige,
mellem Virkelighed af lavere og af høiere Orden;
men lidt mere Varsomhed med at trække de mange
Mirakelhistorier paa Traad og fremstille det Overnaturlige
paa en saa sandselig udvortes Maade, at det første
Skridt paa Erfaringens og Selvtænkningens Vei maa
lede de Unge i en farlig Fristelse, kunde man dog nok
anvende.

Det er ikke Forslag til Forbedringer i det bestaaende
Skolevæsen, vi her fremsætte, saa lidt som
det er Indførelse enten af nye Lærebøger eller nye
Bibelhistorier, vi anbefale. Vi gaae tvertimod ud fra
den Forudsætning, at Skolen, som den nu er, fortjener al
Anerkjendelse --- som Kundskabsskole. Kun dette maae vi
paa det Eftertrykkeligste fremhæve, at saalænge der,
selv hvad Religionsundervisningen angaaer, ikke gives
andre Folkeskoler end de efter det lærde Forbillede indrettede
Kundskabsskoler: da kunne disse Kundskabsskoler,
de være iøvrigt saa fortræffelige som de ville,
dog aldrig grundlægge enten sand personlig Dannelse
eller »levende Folkeoplysning«. Ved enhver frugtbar
Undervisning maa enten Kundskabsinteressen eller Livsinteressen
være det Afgjørende; at ville med en vis Ligelighed
gjøre Fyldest i begge Henseender og, som man
siger i det jevne Sprog, slaae to Fluer med eet Smæk,
er at ville det Urimelige, det Modsigende, det Umulige.
% --- p. 322 ---
\setcounter{footnote}{0}%
\opage{322}Saalænge Erhvervelse af Kundskaber er Dannelsens
Maal, og de erhvervede Kundskabers Mangfoldighed Dannelsens
Maalestok, saalænge vil Befolkningens store Flertal
i ethvert Land leve hen i en ufri, almueagtig Tilstand;
thi hvad Kundskabserhvervelse angaaer, kan
Flertallet umuligt konkurrere med et Mindretal af mere
Begunstigede. Men dersom det store Flertal af Mangel paa
Kundskaber ikke er i Stand til at naae en sand, personlig,
selvstændig Aandsdannelse, da vil Fremtiden, trods
politiske Rettigheder og frie Forfatninger, ikke blive
en Rettighedernes, Lighedens og Frihedens Tid, men paa
den ene Side en Kulturforfinelsens, paa den anden
en Raahedens og i begge Henseender en Lovløshedens,
Forvirringens og Trældommens Tid; thi den udvortes
Frihed kan kun begrundes paa den indvortes. »Der gives
og kan gives et frit Folk, men der gives ikke og kan
ikke gives en fri Almue; thi Almuens Trin er, aandelig
seet, Umyndighedens, altsaa Ufrihedens Trin. Det gjælder
ved folkelig Undervisning om at bringe den menige
Mand saavidt, at han vil kunne tænke klart sine egne
Tanker, gjøre sig tilstrækkelig Rede for sine egne Handlinger;
indsee og forstaae, at Ingen i Samfundet kan
have Rettigheder, uden at han tillige har Pligter; indsee
og forstaae, at Enkeltmands Interesse kun fremmes paa
fornuftig Maade derved, at Samfundsinteresserne varetages
i det Store; indsee og forstaae, at Samfundsinteresserne i det Store umulig kunne fremmes, med mindre
de forskjellige Arbeidsklasser, ikke blot de, der
arbeide i det Legemlige, men ogsaa de, der arbeide i
det Aandelige, nyde fuld Anerkjendelse, hver efter de i
Forholdenes Natur begrundede Vilkaar. Kun den, som
% --- p. 323 ---
\setcounter{footnote}{0}%
\opage{323}er i Stand til at give sig selv Grunde for, hvad han vil
og hvad han gjør, vil være i Stand til at forstaae de
Grunde, der gives ham af Andre. Vanskeligheden ved
Menigmands Frigjørelse er ikke den, at han mangler
enten ædel Selvfølelse eller Forstand eller redelig Villie;
men Vanskeligheden er, at hans Tænkning er uklar, hans
Synskreds indskrænket, hans Sind mistænksomt. Vanskeligheden
er, at han ikke ret veed enten at give Grunde
eller tage imod Grunde, fordi han er væsenlig ufri.
Menigmand er ufri, saalænge hans Indre er dunkelt;
men, at hans Indre er dunkelt, giver sig tilkjende derved.   % PRINTED AS IS (p. 323): „derved. at“: a stop before a lower-case word
at han i de mest personlige Anliggender er mørk,
indesluttet. --- Der er i Folket en dobbelt Trang, Trang
til Religion og Trang til Oplysning, Trang til Sjælefrelse
og Trang til Kultur«\footnote[1]{Om Hindringer og Betingelser for det aandelige Liv i Nutiden. S. 287---88.}; skal denne Trang afhjælpes,
da maa der begyndes forfra paa en ny Vei,
med en ny Undervisning, thi Aarhundreders, ja Aartusinders,
Erfaring har lært, at Maalet ikke naaes med
den gamle.
\signature{R. \emph{Nielsen}.}

\clearpage
\essayhead{To Slags Undervisning.}{II.}
% --- p. 409 ---
\setcounter{footnote}{0}%
\opage{409}Dersom Forskjellen mellem de nye og de gamle Skoler
alene bestod deri, at hine ligefrem tilbøde en Fortsættelse
og Udvidelse af den Undervisning, der allerede var
grundlagt i disse, da vilde det forhaabede Fremskridt
vist ikke have stort at betyde. Skulde der uden videre
bygges paa det gamle Grundlag, da maatte man jo ogsaa
i alt Væsenligt anvende den gamle Fremgangsmaade.
Om der ogsaa valgtes et nyt Stof, idet man i
Bondehøiskolerne særlig lagde an paa at gjøre Menigmand
bekjendt med Fædrelandets Sagnhistorie, Literatur
og politiske Tilstande, det vilde dog ikke forandre
Sagen. Thi hvorledes skulde man indrette Fremgangsmaaden?
At gjøre Kundskabsmeddelelsen lidt friere,
livligere og fornøieligere ved udelukkende Brug af det
mundlige Ord og ved at fritage Eleverne for Lektielæsning
og Examen, skulde neppe føre til Maalet; Følgen af et
saadant Frihaandsforsøg vilde vel snarere blive, at det
Nye forstyrrede det Gamle uden selv at komme til Kræfter.
Nei, Sandheden er, at der fra Kundskabsskolen til
Folkeskolen ikke gives nogen ligefrem brolagt Overgang.

Rodfæstet i det Gamle er den doktrinære Undervisningsmand
sat i den pinlige Forlegenhed, at han ikke
% --- p. 410 ---
\setcounter{footnote}{0}%
\opage{410}kan Andet end føle den Livsmagt, der aander i det Nye, og
dog maa stirre paa dette Nye med et ærgerligt og forarget
Blik, et Blik, der standser ved Skallen uden at
opdage Kjernen. Fra denne Kant ville adskillige Indsigelser
da ogsaa naturlig lade sig høre. »Man begynder
nu med at gjøre Folkeundervisning til en Art behagelig
Tidsfordriv, en Slags Ordsprogsleg med Sang og
Klang og alskens opbyggeligt Morskab. Men at arbeide
aandelig er ikke det Samme som at more sig. Det er
»Alvor og Gammen, der skal kunne saa godt tilsammen«;
og saa vælger man den Gammen, der passer sig,
ikke for Læreren, men for Ammen. Det er jo at nedværdige
den guddommelige Fornuft. Al sand Undervisning
maa begynde med Ærbødighed for det, som er
høiere end Mennesket, saa at Enhver, der træder ind i
en stor Læreanstalt. kan mindes, hvad Ordet siger:   % PRINTED AS IS (p. 410): „Læreanstalt. kan“: a stop before a lower-case word
drag dine Sko af, thi dette Sted er helligt. Der er en
ammestueagtig Naturlighed, hvormed man frit efter Lune
og Indfald tør give sig i Snak om Livets allerhøieste
Anliggender, og kige »Vorherre« kammeratlig i Kortet:
mon det skulde være det Nye? I den selvbehagelige
Indbildning, at djærv Smagløshed og logisk Ugenerthed
afgive det fortræffeligste Bevis paa folkelig Veltalenhed,
fylder man Hovederne paa de unge Lærlinge med Viser
og Ordsprog, Emter og Eventyr, folkehistoriske og verdenshistoriske   % PRINTED AS IS (p. 410): „Emter“
Overblik saalænge, at det tilsidst maa
løbe rundt for Øinene. Ifølge en mere drømmende end
dømmende Opfattelse af Aandsfænomener i Billedsproget
forklarer man Billede ved Billede, og nikker beundrende
til ethvert Mundheld uden at kjende en Tanke.
Langtfra at gjælde for hvad de oprindelig ere, Symbo%
% --- p. 411 ---
\setcounter{footnote}{0}%
\opage{411}ler for Oldtidstanker, som kun i en given Sammenhæng
have en vis indskrænket Almengyldighed, blive
Sagn og Myther behandlede som Spaadomme af uudgrundelig
dybt Indhold, og Fortolkningen gjort afhængig af
den allervilkaarligste Ide-Association. Hvad »Mythesmeden«
har ment, er ligegyldigt, det Vigtige er, at
Mythen kan bruges til hvadsomhelst: en Billedæske,
hvori man gjemmer, hvad Skjaldeblikket har seet, et
koloreret Omslag om den profetiske Aands allernyeste
historiske Aabenbaringer. Det er en falsk Stræben efter
at være paa en Gang dybsindig, aandrig og dog enfoldig.
For at modarbeide en Yderlighed kaster man
sig over i en anden. Man vil undgaae de stive Former
og bliver formløs; man vil bort fra det Konstlede og
forfalder til Plathed; man forsmaaer forfængelig Pynt,
og griber det Rene med uvaskede Hænder. En saadan
Ligefremhed, hvorved man paa snaksom Maade familiariserer
sig med Tilværelsens ideale Magter, gjør et Indtryk
omtrent, som hvis en Menighed i en af Fristaterne,
for at undgaae det falsk Høitidelige, faldt paa at samle
sig og holde Bøn med Hænderne i Lommen og Cigar
i Munden.« --- Lad dette saa være en Prøve paa doktrinær
Forstaaelse; det er jo ganske naturligt, at de,
der ikke formaae at trænge ind i Fæstningen selv, maae
holde sig til Udenværkerne.

Grundforskjellen mellem det Nye og det Gamle er
iøinefaldende, naar vi see hen til de første Betingelser.
Den gamle, officielle Folkeskole er en Børneskole. Mod
Børn og Umyndige kan der anvendes Tvang; Børneskolen
under Statens Formynderskab er da ogsaa en Tvangsskole
i dobbelt Henseende: for Børnene og for Foræl%
% --- p. 412 ---
\setcounter{footnote}{0}%
\opage{412}drene. Den nye Folkeskole er en Skole for Voxne;
for de Voxnes Undervisning har Samfundet ingen Tvangslove:
her kan altsaa begyndes med Frihed. Den gamle
Folkeskole er en Kundskabsskole og maa som Kundskabsskole
bygge paa et konstigt, mekanisk Grundlag;
den nye Folkeskole derimod, som ved Vækkelse af de
personlige Aandskræfter vil være en Skole for Livet,
kan tale med de levende Mennesker om Livet selv, om
hvad der ligger dem nær og begynde med et naturligt
Grundlag. Ved at betragte disse Udgangspunkter
hvert for sig lidt nærmere ville vi let overbevise os om,
hvor væsenlige de ere.

Det er en Omstændighed, hvorpaa der maa lægges
Vægt, at Ungdommens officielle Undervisning netop ender
i det Tidspunkt, hvor Aandslivet egenlig først for
Alvor skal til at begynde. At man ikke ved Tankelektier
til Udenadslæsning kan foregribe den modnere Alders
forventede Eftertanke, er allerede viist. Det gaaer
ikke med Levevisdom som med en Konfirmationskjole, der
kan syes for lang og vid, for at blive tilpas med Aarene.
Ved at give Barnet et systematisk Udtog af
Sædelæren fremmer man ikke dets Sædelighed; endnu
mindre styrker man dets Tro paa Gud ved en i Lærebogsform
affattet Børnedogmatik, thi Børnedogmatiker
og dogmatiske Børn skulle vanskelig »mødes i Guds Rige«.
Hvad der gjælder om Religion og Sædelære, gjælder paa
tilsvarende Maade om Verdenshistorie, Fædrelandshistorie,
det borgerlige Livs Fordringer osv., der er en Grændse
for de Umyndiges Fatteevne, som ikke kan overskrides,
uden at der gjøres Brud paa den guddommelige Orden.
Det er i denne Henseende af Vigtighed, at den nye
% --- p. 413 ---
\setcounter{footnote}{0}%
\opage{413}Folkeskole begynder Undervisningen forfra, idet den begynder
med de Voxne.

Er det allerede i sig selv unaturligt, at en Skole for
Voxne skulde nedlade sig til at virke ved Tvang, saa
bliver det Unaturlige dog endnu mere unaturligt, naar
det udtrykkelig er en Vækkelsesskole, en Skole for Livet.
Ved den blotte Kundskabserhvervelse, navnlig ved
Alt, hvad der gaaer ind under Lektielæsning og Examen,
er det en temmelig ligegyldig Sag, om Indøvelsen
er skeet frivillig eller ufrivillig. Naar en Dreng har lært
den store Tabel, saa har han lært den, og saa er det
Pust overstaaet. Om det gik med det Gode eller med
det Onde, om det var ham en sand Fornøielse eller det
kostede Taarer, kan han gjerne glemme, naar han blot
ikke glemmer Tabellen. Anderledes med de Sandheder,
der skulle virke paa Sindet og have Betydning for Personligheden.
Vistnok gives der i det praktiske Liv
mangfoldige Handlinger, som hverken ere aldeles frivillige
eller aldeles ufrivillige, men en Blanding af begge,
saasom at betale Skat, at gjøre Krigstjeneste o. desl.
Tvang og Frihed kunne forliges i alle saadanne Tilfælde,
hvor det mere kommer an paa Handlingens Udvortes
end paa dens Indvortes, mere paa Gjerningen end paa
Bevæggrunden, i det Hele mere paa Samfundet end paa
Individerne. Men Livsanskuelse, personlig Overbevisning,
Kjærlighed til Fædrelandet, levende Kjærlighed til Gud
og Næsten, slige Aandens Gaver lade sig ikke inddrive
ved Tvang, de ere Frihedens Frugter, og Friheden har
sin Begyndelse i sig selv.

Det Begreb, hvori Enhed af Tvang og Frihed oprindelig
maa søges, er Begrebet: Opdragelse. Opdragelse
% --- p. 414 ---
\setcounter{footnote}{0}%
\opage{414}forudsætter Optugtelse og er forsaavidt betinget af Tvang;
men sand Opdragelse maa altid ende med Selvopdragelse.
Optugtelse ved Andre faaer kun Betydning ved
at lede til Selvoptugtelse; Selvoptugtelse kræver en fast
Villie, en kraftig Selvbestemmelse: her begynder Friheden.
De to i Opdragelsen liggende Opgaver kunne nu
ligelig fordeles paa Modsætningen mellem den gamle og
den nye Skole.

Løsenet i det gamle Opdragelsessystem har altid
været, og vil forhaabenlig vedblive at være: Respekt og
Lydighed. Med denne imponerende Fordring træder
Kundskabsskolen de Unge imøde, og benytter saa de Betingelser
og Midler, Undervisningen stiller til dens Raadighed,
for at lade dem føle, at der er Noget til, der hedder
Lov. At være i Skolen er at være under Loven.
Den strenge Disciplin er ikke et nødvendigt Onde, nei
den er et væsenligt Gode, thi fornuftig Strenghed er vidt
forskjellig fra despotisk Haardhed. Strenghed er Alvorens
Udtryk, der gjør det Ubetingede føleligt og rammer
Uvorenheden med Digterens Ord:

\begin{quote}
»Respekt er hvad det skorter paa,\\
Indprentes maa den i Din Hu,\\
Og Du Dig underordne maa\\
For det, som er lidt Mer end Du.\\
Respekt, som sagt, thi nu har Du ingen,\\
Respekt for Stedet, Respekt for Tingen.«
\end{quote}

Det er vel fornemlig den Respekt for Tingen, hvortil den
gamle Skole havde optugtet sine Disciple, som Slutningen
af forrige og Begyndelsen af dette Aarhundrede
endnu kunde takke for saamange store offenlige, handlekraftige
Karakterer. Karakter forudsætter fast Overbe%
% --- p. 415 ---
\setcounter{footnote}{0}%
\opage{415}visning og stærk Villie; men fast Overbevisning og
stærk Villie forudsætter igjen, at der er Noget, der
ubetinget vil respekteres. Ved Respekt for Tingen kunde
da Striden mellem modsatte Livsanskuelser slaae igjennem
i Samtiden og blive lidt Mere end et Opinionskævleri
mellem avisskrivende Individer: hver Anskuelse, der vilde
staae og falde med sin Karakter, blev agtet; kun den
vindige, Intet villende Karakterløshed var foragtet.

Hvor Respekt for Tingen mangler, indretter man sig
paa en anden Maade. I den Forudsætning, at enhver
bestemt Anskuelse, ligesom ogsaa enhver bestemt Karakter,
er en Ensidighed, en Art Bornerthed i Villien,
indretter man sig en literær Butik med Principer og
Theorier i alle Mønstre, gamle og nye, positive og negative,
frem- og tilbagegaaende efter Omstændighederne, et Kaleidoskop
af de mest modstridende Anskuelser. Hver Gang
nu En viser sig med en ny Anskuelse, er der altid en
Svend i Butiken, der passer ham op, idet han gaaer
forbi, kaster ham til Publikums Moro en gammel og
modsat Anskuelse lige i Øinene og griner ham ud, ikke
fordi den Forbigaaendes Anskuelse just er daarligere end
saamange Andres, men fordi han har den Anmasselse,
den Villiens Bornerthed, at ville have en bestemt Anskuelse
og ligne Noget, der kunde kaldes Karakter.
Lykkes det paa denne Maade lidt efter lidt at forvandle
et høistæret Publikums Respekt for Tingen til Respekt
for Butiken, saa er Spillet gaaende. Saa kan en
hæderlig Eier af en saadan hæderlig Butik sidde lunt
i Offenlighed og slaae sig op i al mulig Hæderlighed;
det skulde da være, at der dog var een Hæderlighed,
han ikke kunde naae, at han ikke kunde naae at blive
% --- p. 416 ---
\setcounter{footnote}{0}%
\opage{416}hæderlig Karakter. Karakter er nemlig ifølge Forudsætningen
stemplet som Ensidighed; men stemplet Ensidighed
sprænger Butikken, thi Butikken er den luftige
karakterløse Alsidighed. At nu et Individ, der ikke kan
blive Karakter, heller ikke kan blive hæderlig Karakter,
er da vel ligesaa indlysende, som at et Væsen, der
ikke kan blive Menneske, heller ikke kan blive Medmenneske.
»Respekt er hvad det skorter paa«, Respekt,
ikke for Individet, men for Sagen, naar Sagen er et
Livsspørgsmaal, hvis Besvarelse umulig kan undgaaes,
medmindre Karakterløsheden selv skal være Tidsalderens
Karakter. Mangel paa Respekt for en vigtig Sag
er Mangel paa Opdragelse, et Vidnesbyrd om, at de
ræsonnerende Individer have unddraget sig den strenge
Disciplin ved at løbe for tidlig af Skole.

Den strenge Disciplin er den gamle Skoles Hæder;
den staaer paa Lovens Grund og peger mod Frihedens
forjettede Land, men kommer ikke selv derind; hvor
Tvangen ophører, har den sin Grændse. Ved denne
Grændse tager den Afsked med Disciplene, som nu paa
egen Haand maa vove det leukadiske Spring ind i den
nye Region. Overgangen fra det Gamle til det Nye er
betegnet ved det Vendepunkt, da Tvangsopdragelsen
skal forvandles til Selvopdragelse. Med Hensyn paa den
høiere Folkeskole er denne Overgang udtrykkelig gjort
afhængig af et Valg. Ikke hvad der vælges, men at
der vælges, er her det Afgjørende. Saa langt som
Tvangen rækker, falder Ansvaret væsenlig paa Formyndere
og Lærere; naar Tvangen ophører, falder Ansvaret
paa den Paagjældende selv. Det bliver nu hans
egen Sag, om han vil nøies med, hvad han har, eller
% --- p. 417 ---
\setcounter{footnote}{0}%
\opage{417}tragte efter noget Bedre. Frihedsskolen indbyder ham,
tilbyder ham sin Tjeneste, men den frister ham ikke.
Den har ingen endelige Fordele at tilbyde ham, ingen
lokkende Udsigter, hvad det Udvortes angaaer; den
maa tvertimod begynde med at fordre endog forholdsvis
betydelige Offere af Penge, Tid og Arbeidskræfter.
At Saamange endog blandt de Ubemidlede allerede
have modtaget Vilkaarene, er et glædeligt Tidens
Tegn. Det vidner om, hvad der rører sig dybt i Folkets
Sind: en Trang til Vækkelse, til indre Frigjørelse.
Det menige Folk aner, hvad fornem Klogskab
ikke har kunnet udregne, at der maa gives en Sandhed,
som uafhængig af endelig Forskjel i Kundskab og Lærdom
er tilgængelig for Alle, en levende Sandhed, som
tænder sin Fakkel ved Frihedens Sol og gjør aandeligt
Formynderskab overflødigt. Der er da al psykologisk
Sikkerhed for, at unge Mænd, som under de angivne
Betingelser indtræde i Frihedsskolen, ikke ville enten
misforstaae eller misbruge Friheden. Den høiere Kundskabsskole
maa altid spørge de Indtrædende om Forkundskaberne,
thi i Forkundskaberne bestaaer Forberedelsen;
Frihedsskolen derimod spørger ikke om Forkundskaberne,
den sætter Forberedelsen i Trangen. Ikke
den Elev, der har de fleste Kundskaber, men den, der
har den mest opvakte Sands og føler den stærkeste Tilskyndelse
til virksom Deltagelse i det aandelige Arbeide,
er den bedst Forberedte. Han kommer da heller ikke
for et Skins Skyld, dertil er Opoffrelsen for kostbar;
nei, han kommer for Sagens Skyld. Han vil ikke døsig
gnave paa en Tankelektie, ubekymret om Tanken
selv, kun opmærksom paa, hvordan Læreren »vil have
% --- p. 418 ---
\setcounter{footnote}{0}%
\opage{418}det«; han vil, saavidt det staaer i hans Magt, anstrenge
sig for at see med egne Øine, høre med egne Øren,
tænke med egen Hjerne: en saadan frisk Oprindelighed er
vel i Stand til at opveie Savnet af mangeaarige Forstudier.

Grundlaget for Undervisningen giver sig nu af sig
selv; thi Skolen begynder med Begynderen der, hvor han
er, og taler kun til ham om Ting, han uden Vanskelighed
strax kan forstaae. Til at tyde et Ordsprog eller
høre et Folkesagn, et Eventyr, et Stykke Historie er
Enhver forberedt; den daglige Erfaring har desuden beriget
ham med et ikke ringe Kjendskab til Menneskelivet
og den ham omgivende Natur: her er Tilknytningspunkter
nok, naar Læreren forstaaer at benytte dem.
Men hvorledes skulle de benyttes? Her møder en Vanskelighed,
vi ikke tør oversee, og det saa meget mindre,
som den let kan forstørres af dem, der ønske at
gjøre Indvendiger. »Lad Begyndelsen« --- siger man ---   % PRINTED AS IS (p. 418): „Indvendiger“ for Indvendinger
»være nok saa naturlig, og det mundlige Ord nok saa
vækkende: saa er her jo dog et Stof. Stoffet vil voxe
hver Dag med hver Time: hvorledes skal nu Tilhøreren
ved den blotte Høren kunne samle og ordne det ophobede
Stof, dersom han hverken har Bøger eller skriftlige
Optegnelser? Men dersom han for at have Gavn
af Foredraget maa tye til Lærebøger og skrevne Hefter,
saa bliver jo Fremgangsmaaden ved Folkehøiskolen i alt
Væsenligt den samme som ved den lærde Høiskole. For
at det mundlige Ord skal kunne bære sin Frugt, maa
Eleven i Forveien have lært baade at læse og skrive;
han maa begynde med Kundskaber for at erhverve sig
Kundskaber: paa denne Maade bliver Skolen for Livet jo
dog en Kundskabsskole; og hvad Forskjel er her saa
% --- p. 419 ---
\setcounter{footnote}{0}%
\opage{419}imellem det Nye og det Gamle? At Folkehøiskolen ikke
har kunnet holde sig paa Høidepunktet af bare Vækkelse
og Begeistring, men har maattet nedlade sig til det
Reale og fra Begyndelsen indrette sig paa at meddele
nyttige Kundskaber i bestemte Fag, er en Kjendsgjerning.
Saaledes finder man f. Ex. i »Beretning om Folkehøiskolen
i Rødding af C. Flor 1846« Timetabeller,
der lyde paa »Theoretisk Landvæsen«, »Fysik«, »Geografi«,
Statistik«, »Arithmetik og Regning«, »Geometri
og Tegning« osv. --- Skulde den stærkt betonede Forskjel
mellem det Gamle og det Nye have principiel
Gyldighed, da maatte Fremskridtet jo være betegnet ved
Opdagelsen af en ny Methode, en Undervisningsmaade,
der gjorde det muligt, at bibringe uforberedte Elever
grundige Kundskaber uden at uleilige dem med Lektielæsning
og Examen. Er en saadan Opdagelse virkelig
gjort? Og hvis den var gjort, hvad saa? Saa vilde
den jo strax komme den gamle Skole tilgode og netop
tjene til at udjevne enhver Kløft imellem det Gamle og
det Nye.«

At Indvendinger af denne Natur beroe paa en aabenbar
Misforstaaelse, er let at vise. For at Kundskabsskolen
kan svare til sit Begreb, maa Kundskabsmeddelelsen
være Maal, ikke Middel; i Folkeskolen derimod
maa Kundskabsmeddelelsen være Middel, ikke Maal. Det
gjør en væsenlig Forskjel. Talen være f. Ex. om Geometri.
Dersom der nu begyndes saaledes: Geometri er
Videnskaben om Rumstørrelser, et System af sammenhængende
Læresætninger; det er desaarsag nødvendigt at skride
frem fra Bevis til Bevis uden at overspringe noget væsenligt
Mellemled; ethvert Bevis faaer sin Betydning ved
% --- p. 420 ---
\setcounter{footnote}{0}%
\opage{420}de Forudsætninger, hvoraf det fremgaaer, og de Følgesætninger,
hvortil det videre kan føre osv., saa høre vi
strax, hvor vi ere; vi ere i Kundskabsskolen. I Folkeskolen
begyndes der fra et andet Udgangspunkt. Det er
ikke Bevisførelsen i og for sig, ikke Relationerne mellem
mathematiske Linier, Vinkler og Flader, der oprindelig
vækker Interessen; det er i Reglen Opgaver af rent
praktisk Art. At en Trekant altid er fuldkommen bestemt,
naar to Vinkler med en hosliggende Side ere
givne, vil neppe i nogen synderlig Grad tiltale Begynderen.
»Men hvorledes er det muligt, at man staaende
paa Jorden kan maale Høiden af et Taarn eller af et
Træ i Skoven? Den Opgave var det nok værd at faae
løst«. Ved at gaae ud fra noget Virkeligt, fra Noget, der
er ham umiddelbart bekjendt, ledes Eleven let til en
Forstaaelse af, hvad Læren om Forholdet mellem Trekantens
Dele vil sige. Interessen for Bevisførelsen afledes
af Interessen for den praktiske Opgave. Det Antal
Læresætninger og Beviser, han maa gjennemgaae, og den
Orden, hvori disse følge paa hinanden, retter sig altsaa
ikke efter Systemet, men efter den valgte Opgave. Det
er end ikke nødvendigt, at alle Læresætninger blive beviste.
Naar Eleven først har begyndt at indsee, hvad
et Bevis er, kunne mange Sætninger ligefrem tages paa
Autoritet. Ved et Gilde paa Landet, hvor Samtalen i
Anledning af Almanakken havde begyndt at tage et
høiere Sving, var der engang En, der spurgte: »hvor langt
monstro er der fra Jorden til Maanen«? »Det er der ingen
Mennesker, der veed«, svarede en gammel Bonde; »der er
endnu Ingen af de Lærde, der har kjørt den Vei.« --- Skulde
en Elev uden Forkundskaber arbeide sig systematisk
% --- p. 421 ---
\setcounter{footnote}{0}%
\opage{421}igjennem Arithmetiken og Geometrien ind i Astronomien,
indtil han i Kapitlet om Parallaxen kunde faae hint Spørgsmaal
besvaret, da vilde Arbeidet i Forhold til den korte
Tid, han kan anvende derpaa, ganske vist være uoverkommeligt.
Og dog kan Besvarelsen af et eneste saadant
Spørgsmaal være af overordenlig Betydning for hans
personlige Dannelse; den kan aabne ham et Indblik, ikke
blot i Verdensbygningen men ogsaa i den menneskelige
Viden, et Indblik, som hæver ham høit over det almueagtige
Standpunkt og giver ham en Forestilling om, at
Videnskab og almindelig Oplysning have Noget med
hinanden at gjøre. Heldigvis kan han, hvis det virkelig
er ham om at gjøre, med en let overkommelig Forberedelse
faae sit Spørgsmaal besvaret, uden at segne
under Systematikens Byrde. Det Samme gjælder om
de Hovedbeviser, paa hvilke det Kopernikanske System
støtter sig. Her kan ved praktisk Fremgangsmaade erhverves
Kundskaber, endog sammenhængende Kundskaber,
f. Ex. i Landmaaling; men Kundskaberne ere Midler,
ikke Øiemed. Det Omfang og den Sammenhæng,
hvori de meddeles, er afhængig af Elevens individuelle
Begavelse og personlige Interesse. Det er ikke for Videnskabens
Skyld, men for sin egen Skyld, at han arbeider;
det er ikke Mængden af det, han lærer, men
den Veiledning, han faaer til at udvikle sig selv, der er
Hovedsagen; det er Vækkelsen, hvorpaa det først og
sidst kommer an. »Det Vigtigste ved vor Høiskole« ---
hedder det i anførte »Beretning« S. 13 --- »er ikke de
positive Kundskaber og Færdigheder, som vi søge at
bibringe Eleverne, men snarere det hele aandelige Liv,
som bliver vakt og næret hos dem, saa at deres For%
% --- p. 422 ---
\setcounter{footnote}{0}%
\opage{422}stand bliver skarpere, deres Dömmekraft modnere, deres   % PRINTED AS IS (p. 422): „Dömmekraft“, set with ö
Hjerter aabnere og ædlere, at der vækkes hos dem
Sands for Orden, Skjønhed og et smagfuldt Liv, at Lyst
til Beskæftigelse kan træde istedenfor den sædvanlige
Dorskhedsvane, at deres Sind og derved deres Væsen
kan blive frimodigere, deres Følelse for Kammeratskab
og Fædreland blive vakt, næret og styrket«.

Gjælder det allerede om den Kreds af Kundskaber,
der nærmest angaaer de udvortes og endelige Ting, at
den afgiver Vækkelsesmidler for personlig Selvudvikling,
da gjælder dette endnu i høiere Forstand om de Dannelseselementer,
der ifølge deres Natur udelukkende tilhøre
Aandens Rige; vi tænke paa Mythologi og Historie
og begges Forhold til den aabenbarede Religion.

»Hvem der vil have evig Krig« --- siger N. F. S.
Grundtvig i sin Fremstilling af »Nordens Kæmpeaand«
(Nordens Mythologi 1832, S. 95 flgd.) --- »har nemt ved
at naae sin afsindige Hensigt, thi han behøver blot at
forkynde Ærlighed og Sandhed aaben Feide, da skeer
hans Villie; men hvem der vil have evig Fred, og hvem
skulde da ikke altid ville det? han slipper ingenlunde
saa nemt dertil, som de kjære Filosofer allevegne har
bildt dem ind, og som man nu sædvanlig tænker. Fred og
Ro er nemlig Et, og Stilhed er et Andet, saa Rolighed
lader sig forene med stærk Bevægelse, og Stilhed kan
være forbundet med stor Uro, som man vel klarest seer
ved at sammenligne den ædelmodige, seirrige Helt med
den bagbundne Slave. Dette var ogsaa Billedet, der
svævede for vore fribaarne, høimodige Stamfædre, saa
de følte dybt, at Livet er den første Betingelse for hvad
levende Folk kan kalde Fred og Ro, og Kræfternes frie
% --- p. 423 ---
\setcounter{footnote}{0}%
\opage{423}Udvikling den anden\dots{}. Denne nordiske Anskuelse af
det timelige Liv hos Guder og Mennesker som en bestandig
Kamp, det kun gjælder om at føre saa ædelt,
saa kraftigt og saa klogt, som muligt, da al den virkelige
Fred, der kan times Gudefolket og Adelstammen i
Tidens Løb, udspringer som en Seirsgudinde af Kampen;
denne Anskuelse maa i en vis Grad findes hos
ethvert Folk, der skal hævde sin Selvstændighed og
faae Lov til at udvikle sine Kræfter, selv paa den fredeligste
Maade, og at den eiheller i Norden er uddød, har
det nærværende Aarhundrede viist; men sit aandige Præg   % PRINTED AS IS (p. 423): „aandige“ (so also p. 618)
har den dog næsten tabt, saa, istedenfor at udvikles til
Klarhed, fordunkledes den mer og mer, til det blev
den almindelige Mening, at Krige vel er et nødvendigt
Onde, fordi der gives rovgjerrige Folk, ærgjerrige Kæmper
og herskesyge Fyrster; men at al Kamp dog i
Grunden er afskyelig, og mer afskyelig, jo aandeligere
dens Gjenstand er\dots{}. Man kan vel tænke, at det er
ei min Hensigt dermed at forgude Asa-Thor eller saaledes
at vaske Blodet af Nordens Kæmper, at de synes
englerene, thi alle Historiens Kyndige kan see, at ligesom
der kun har været een Helt med uplettet Skjold,
saaledes er ogsaa Kristendommens Kæmpe-Aand den
Eneste, der fortjener Guddomsnavn og Tilbedelse; men
vel ønskede jeg ret at indskærpe alle mine Læsere den
miskjendte Sandhed, at, uagtet Kæmpeaandens jordiske
Redskaber var ligesaa langt fra at være lydesløse, som
Fredsmæglerne, saa var de dog derfor hverken mindre
uundværlige eller mindre berettigede til Beundring og
skaansom Bedømmelse i en syndig Verden.«

Med disse Ord er den afgjørende Synsretning, saa%
% --- p. 424 ---
\setcounter{footnote}{0}%
\opage{424}vidt vi kunne skjønne, engang for alle angivet: Mythologi,
Historie og Religion belyse hverandre gjensidig.
--- Man kan være uenig om Maaden, paa hvilken den
udtalte Livstanke bedst lader sig udvikle og sætte i
Kraft, men om Tanken selv maae Alle, der ville lidt
Mere end en nysgjerrig Kigen i Døgnkombinationernes
Kaleidoskop, upaatvivlelig være enige. Forskjellen mellem
Kundskabsskolen og Folkeskolen er her saa slaaende,
som muligt. Thi naar Eftertrykket ensidig lægges paa
Kundskaberne, maae Mythologi, Historie og Religion naturligvis
falde ud fra hverandre i tre Fag, tre særskilte Stykker,
der have Lidet eller Intet med hverandre at gjøre.

Af Mythologien som Kundskabsfag lærer man da en
hel Del om nordiske, græske, romerske Guder med tilhørende
Gudinder: det er mange Navne, mange Figurer,
Attributer, Optrin og eventyrlige Bedrifter, Eleven har
at indprente sig, naar han skal finde Rede i det Altsammen.
Og hvad kommer saa alt dette os ved; vi troe
jo dog ikke længer paa de gamle Guder? »Vistnok« --- svarer
man --- »hvad Frihedslivet, vort indre sjælelige Liv
med den nugjældende Verdensanskuelse, angaaer, kommer
det Hele os egenlig slet ikke ved: vi ere jo langt
udover hine gamle Fantasier«. Og hvorfor skal man
saa lægge sig efter at kjende dem? »Det er nødvendigt
for Dannelsens Skyld; det er ganske nødvendigt at kjende
Noget til Oldtidens ideale Forestillinger. De mythiske
Digtninger, i hvilke det Ophøiede og Rædselsfulde ofte
forbinder sig med Træk af barnlig Ynde og naiv Elskværdighed,
indeholde desuden en rig Mangfoldighed af
kunstneriske og poetiske Motiver. Hvorledes skulde man
kunne forstaae de mange Allusioner hos Nutidens Digtere,
% --- p. 425 ---
\setcounter{footnote}{0}%
\opage{425}dersom man ikke var nogenlunde orienteret i sin Mythologi.
Den, der ikke veed, at Frigga var Odins Hustru,
at Venus var gift med Vulkan, at Mars var Krigsgud
hos Romerne ligesom Thor hos Nordboerne, maa jo
sidde som en Idiot i dannede Menneskers Selskab, kan
ikke besøge et Musæum, ikke engang gaae paa Udstillingen
og lukke sin Mund op uden at prostituere sig.«
Det er altsaa ikke Mythens Væsen, men kun dens Skin,
et æsthetisk reflekteret Efterskin, for hvis Skyld Anstrængelserne
gjøres.

Og nu Historien? Man har, som bekjendt, Kompendier
af Konsthistorien, der give en Udsigt over de
vigtigste Skoler i Konstens forskjellige Perioder med
Navne paa berømte Konstnere i hver Periode samt en
kort Beretning om hver Konstners Levnet, Arbeider,
Fødeby, Dødsaar osv. Dersom nu et ungt Menneske,
om ikke just af særdeles Kjærlighed til Konsten, saa
dog fordi det nu engang hører med til den høiere Dannelse,
ønskede at sætte sig noget ind i Konsthistorien
og af den Grund gjorde en Reise til en af de Byer, der
udmærke sig ved rige Konstskatte og berømte Gallerier:
da kunde dette jo meget vel lade sig høre. Men dersom
han saa, efter at være ankommen til et Konstens
Hovedsæde, istedenfor at gaae omkring i Gallerierne og
betragte Mesterværkerne, blev siddende i sit Logi for
at læse over paa Kompendiet, indtil han kunde det udenad,
og saa, idet han var lige nær ved at kunne det og derved
omsider forberedt til at nyde Beskuelsen, opdagede, at
hans Tid var udløben, at han nu paa Grund af andre
vigtigere Forretninger maatte forlade Byen, saa vilde man
dog vel finde, at han havde taget Sagen noget bagvendt.

% --- p. 426 ---
\setcounter{footnote}{0}%
\opage{426}Hvad der ved sig selv er indlysende med Hensyn
paa Konsthistorien, synes ikke endnu at være ganske
indlysende med Hensyn paa Verdenshistorien, thi her er
det Kompendiet først og Kompendiet sidst, hvorom Undervisningen
dreier sig. Det koster i Reglen saamange
Aars Anstrengelse at lære et med Navne og Aarstal vel
spækket Udtog af Verdenshistorien nogenledes udenad,
at naar Maalet er naaet, saa er Tiden omme; saa indtræder
den omvendte Bevægelse: man glemmer med Lethed,
hvad man lærte møisommeligt, og det taagede Omrids,
der bliver tilbage, er endnu længere fra at gjøre
Indtryk af virkelig Historie, end Elefanten i en ABC-Bog
er fra at gjøre Indtryk af virkelig Elefant. Et
Riges Erobring kan fortælles i to Linier, en hel Revolution
paa en halv Side; selv om der anvendes et Blad
eller to til Beskrivelse af en Verdensomvæltning, kommer
saa Læseren derved virkelig ind i de virkelige
Begivenheder? Man holder sig til Kjendsgjerningerne
for at gribe det Virkelige, og mærker ikke, at Virkeligheden
selv med dens Strid og Smerte, dens Kræfter
og Lidenskaber er dunstet bort. Hvad lære vi saa af
Historien? »Mange Ting, mærkværdige Ting, overordenlig
Meget!« Og dog, naar Spørgsmaalet er om det
nærværende Øieblik, overordenlig Lidet. Saalidt som
de Døde kunne skrive Love for de Levende, ligesaalidt
kan den forbigangne Tid umiddelbart bestemme, hvad
der skal skee i den nærværende. Vel er det sandt, at
Eftertiden altid staaer paa Fortidens Skuldre; men det
er derfor ikke mindre sandt, at enhver Tidsalder med
Aand og Liv har Opgaver, med hvis Løsning den maa
begynde forfra. »Hvad maa der gjøres?« »Hvad vil
% --- p. 427 ---
\setcounter{footnote}{0}%
\opage{427}der skee?« Naar disse Spørgsmaal blive paatrængende,
da nytter det kun lidet at slaae op i et verdenshistorisk
Kompendium og spørge Fortiden tilraads; thi Fortidens
Orakler ere dunkle. Det er med Fortiden, naar man
paa udvortes, endelig Maade tyer til den om Raad, som
med Spaakvinden i Endor: det Høieste, den formaaer,
er at mane et Gjenfærd. Ethvert dristigt Foretagende
har i Fortidsbegivenhederne altid ligesaamange fraraadende
som tilraadende Grunde; vil den Modige vise de
afskrækkende Exempler fra sig ved at sige: »det var
den Gang, men nu er det anderledes«, da kan den feige
Tvivler jo med de selvsamme Ord afkræfte et hvilketsomhelst
opmuntrende Exempel. Ved ensidig at holde sig
til de saakaldte rene Kjendsgjerninger, de nøgne Fakta,
--- og det gjør Kompendiet paa Kundskabernes Vegne ---
griber man et Udvortes uden dybere Forstaaelse af det
tilsvarende Indvortes. Man mener i et saadant Udvortes
at holde Virkeligheden fast; det er en Skuffelse.
Det udvortes Virkelige i Historiens Kjendsgjerninger er
ikke det Virkelige selv, men kun et Billede deraf, ei
engang et fuldstændigt Billede, kun nogle blege Træk.
Det, som viser sig, er ikke til; men det har været:
fuimus Troes, fuit Ilium et ingens gloria Teucrum. Det
er Skyggen af det Henfarne, der tages for Virkelighed,
og saa bilder man sig ind, at denne Skinvirkelighed er
den historiske Sandhed, den virkelige Historie.

Skal nu Adskillelsen mellem Mythologi og Historie
ligefrem grunde sig derpaa, at Historien giver os det Virkelige,
Mythen det Uvirkelige, og skal saa Intet gjælde
for Virkelighed, uden det, som ligefrem kan bekræftes ved
Vidnesbyrd af de fem Sandser: da vil det være vanske%
% --- p. 428 ---
\setcounter{footnote}{0}%
\opage{428}ligt, for ikke at sige umuligt, at danne sig en sikker
Overbevisning om den aabenbarede Religions Kjendsgjerninger.
Det kan nemlig ikke skjules for et uhildet
Blik, at de bibelske Fortællinger indeholde to uligeartede
Bestanddele, en overnaturlig og en naturlig, af hvilke
den første har en paafaldende udvortes Lighed med det
Mythiske, den sidste med det ligefremt Historiske; og
dog gjøre Aabenbaringsfænomenerne Fordring paa sand
Virkelighed. Det er ifølge Kundskabsudviklingens hele
Gang, især naar vi betænke Naturlærens stedse voxende
Indflydelse, da heller ikke forunderligt, at Kløften mellem
dem, der antage, og dem, der forkaste den positive
Religion, maa blive større og større. Enhver Undergjerning,
hvoraf den hellige Skrift giver en Beskrivelse,
opvækker jo naturlig hos Læseren det Spørgsmaal:
er dette virkeligt eller er det uvirkeligt?

De, som nu, være sig af Mangel paa Dannelse eller
af Mangel paa Eftertanke eller af andre, maaske dybere
personlige, Grunde, aldeles ikke bekymre sig om Naturens
saakaldte »uforanderlige Love«, men uden alle
Mellembestemmelser henføre Alt til Guds almægtige Villie,
kunne jo vel uden Selvmodsigelse antage Undergjerninger
for virkelige; men hvad Modsætningen af det
Naturlige og det Overnaturlige angaaer, vil det, saalænge
de blot holde sig til det Ydre i Kjendsgjerningerne,
være dem umuligt at angive nogen bestemt Grændse
mellem Historisk og Mythisk. Man ræsonnerer i Almindelighed
saaledes: »At Bileams Aseninde virkelig har
kunnet tale, er historisk, thi det staaer i Bibelen, men
at Achilles Hest Xanthos engang skal have havt »det
bevingede Ord«, er mythisk, thi det staaer ikke i Bibe%
% --- p. 429 ---
\setcounter{footnote}{0}%
\opage{429}len, det staaer kun hos Homer«. Et saadant Ræsonnement
kan imidlertid kun gjælde, saalænge det staaer
fast, at Alt, hvad der staaer i Bibelen, er udvortes
bogstavelig historisk. Og dog vil denne Forudsætning,
om hvilken Katholiker og Protestanter officielt ere enige,
ikke engang slaae til. Katholikerne ville, som bekjendt,
ikke indrømme, at Miraklernes Tid er forbi; de antage
tvertimod, at der i Kirkens Historie idelig er skeet
og endnu skeer Mirakler. Den, f. Ex. der ikke kan troe,
at den fromme Kanonikus A er »gaaet tørskoet over
Elben« eller at den hellige Biskop B har kunnet »hænge
sine Handsker paa en Solstraale«, kan heller ikke troe,
at Elias foer til Himlen i en gloende Vogn, og at Jerichos
Mure faldt, fordi Josva lod Præsterne blæse i Basuner;
men den, der trods »den hellige Kirkes enstemmige
Vidnesbyrd« ikke kan troe Sligt, er »en Hedning«,
»en Vantro«, »en Kætter«, »en Tyrk«!

De, som udelukkende holde paa Ligheden, vel at
mærke, den udvortes Lighed, mellem det Mythiske i Hedenskabet
og det Overnaturlige i Kristendommen, kunne,
naar de ville være konsekvente, ikke slaae sig til Ro, før
alt det Overnaturlige er blevet mythisk, og den hele
positive Religion forvandlet til Mythologi.

Der er et vigtigt Punkt, hvorom Atheister og Theologer
altid ere enige, nemlig det udvortes Historiske.
»Er det sandt, hvad Bibelen beretter, saa maa det være
virkelig skeet, er det virkelig skeet, saa maa det være
historisk«, det er den fælles Forudsætning: paa den
udvortes historiske-Virkelighed, skal Sandheden beroe.   % PRINTED AS IS (p. 429): „historiske-Virkelighed“, with a hyphen
Naar Theologien desuagtet som »historisk Videnskab«,
% --- p. 430 ---
\setcounter{footnote}{0}%
\opage{430}vil, hvad den hellige Historie angaaer, komme til en
anden Slutning end Atheismen, da kan den ikke rokke
Oversætningen, men maa forsigtig indskrænke sig til at
lirke ved Undersætningen.

»Dersom det Overnaturlige virkelig er sandt, saa maa
det være historisk.

Nu er det ikke historisk og kan ikke være det.

Ergo er det Overnaturlige i Virkeligheden usandt,
d. v. s. det er mythisk«.

Saaledes er Slutningen atheistisk.

»Dersom det Overnaturlige virkelig er sandt, saa
maa det være historisk.

Nu er det efter al Sandsynlighed virkelig historisk.

Ergo er det i alle Maader vist og sandt, og det
Mythiske udelukket.«

Saaledes er Slutningen theologisk.

Det er det udvortes Historiske, det Skrifthistoriske,
hvorom Apologetiken og den negative Kritik idelig drages
med hinanden. Og saa trækker Apologetiken altid
det korteste Straa; forstaaer sig, det er videnskabeligt
Straa, ellers kunde jo den apologetiske Leg ikke fortsættes.
Den første af de anførte Slutninger er dog idetmindste,
hvad Formen angaaer, logisk korrekt, den
sidste er derimod hinkende, ukorrekt\footnote[1]{Efterat Mænd som G. E. Lessing og S. Kierkegaard tilstrækkelig have godtgjort Umuligheden af „at bygge sin Salighed“ paa objektiv Vished om nogetsomhelst objektivt Historisk, kunde Apologetiken gjerne spare sig Anstrængelsen med at gavne Religionen ved at sammenflikke „Synopser“ og lappe paa „den hellige Historie“.}. Man behøver
blot at besinde sig paa Forskjellen mellem historisk og
% --- p. 431 ---
\setcounter{footnote}{0}%
\opage{431}religiøs Tro, mellem Sandsynlighed og Vished, og man
vil faae en Forestilling om, hvorledes Sagerne staae,
naar Theologien, for at redde nogle historiske Stumper,
maa fire for Kritiken og, for at bevare et Skin af Videnskab,
i Videns Navn nedlade sig til at kokettere
med en let besløret Atheisme.

Ifølge den af Grundtvig angivne Synsmaade, vil Forholdet
mellem Sandhed og Virkelighed vise sig i et
andet Lys. Virkeligheden er dobbelsidig; den har ikke
blot et Udvortes, men ogsaa et Indvortes. Det er
saa langt fra, at Eftertrykket kan falde paa det Udvortes,
at det, hvor Spørgsmaalet er om Sandheden,
meget mere falder paa det Indvortes. Dette gjælder
ligesaa afgjørende for Mythologien som for Historien; thi
det indvortes Virkelige er Aand. Den samme Aand,
der er i et Folks Mythologi, maa ogsaa være i dets Historie,
og omvendt. Historie og Mythologi maa vexelsidig
opklare og belyse hinanden, dersom det ellers skal
være muligt at bestemme Folke-Aanden. Oldtidshistorien
er paa een Gang en Gudernes og en Heltenes Historie:
saadanne Helte, saadanne Guder, og da igjen:
saadanne Guder, saadanne Helte!

Ligesom Heltelivet, det høiere Menneskeliv, har sit Forbillede
i Gudelivet, og kun under denne Forudsætning
kan være virkeligt, saaledes har igjen Gudelivet
sin naturlige Grund og sin Virkeligheds Kraft i Heltelivet.
Forskjellen mellem den græske og den nordiske
Mythologi er i Virkeligheden selv en Forskjel mellem
den græske og den nordiske Folke-Aand, men Forskjellen
mellem den græske og den nordiske Folke-Aand er
% --- p. 432 ---
\setcounter{footnote}{0}%
\opage{432}jo netop den Grundaarsag, hvoraf det Karakteristiske i
begge Folkeslags Historie alene kan afledes\footnote[1]{Forfatteren af denne Afhandling har i Forelæsninger: Om Hindringer og Betingelser (anden Forelæsning) søgt at paavise dette.}.

Det gaaer med et helt Folk som med et enkelt
Menneske: for at samle sig til forhøiet Selvudvikling,
maa det kjende sig selv, sine Kræfter og naturlige Anlæg.
Vidnesbyrdet om dets naturlige Anlæg og deraf
betingede Livsretning er at søge i dets Mythologi; Brugen,
det hidtil har gjort af Kræfter og Anlæg, maa sees
af dets Levnetsløb, dets Historie. Det bedste Middel til
at komme bort fra Øieblikkets Virvar og besinde sig paa
Fremtiden, er at gaae tilbage til Fortiden, men, vel at
mærke, ad Aandens lige Vei, ikke ad de ydre endelige
Begivenheders mangfoldig snoede, ofte vildsomme Stier.
Ved at tabe sig i det udvortes Mangfoldige adspredes
Sindet, Blikket forvilder sig, og Betragtningen ender med
Bedøvelse; kun naar Begivenhedernes Mangfoldighed beherskes
af Tankens Enhed, saa at enhver Udvikling
baade udgaaer fra og leder til Besvarelse af et praktiske
Hovedspørgsmaal, vil Historien bidrage til, at Bevidstheden
klares, at Folket virkelig gaaer i sig selv.
Her er da i Fortiden Meget, der styrker og opløfter,
men ogsaa meget, der ydmyger og nedtrykker, Feil, der
maae rettes, Misgreb og Forsømmelser, der maae angres,
Svagheder, der maae bekæmpes, Skæbnens Slag, der
maae bæres med Resignation. Men for at Historien paa
denne Maade skal kunne virke vækkende og belærende,
maa den ydre Virkelighed overalt henføres til og forklares
af den indre. Historien er da ikke noget blot
Udlevet, et Gravsted for svunden Herlighed, Forbigan%
% --- p. 433 ---
\setcounter{footnote}{0}%
\opage{433}genhedens døde Hav, hvor kun Erindringen boer mellem
sunkne Stæder og Troner, men en nærværende,
bevægende Magt, en Livets Strøm, den Tidernes Elv,
der driver Tilværelsens Hjul.

Først naar det aandige Baand, der forbinder Mythologi
og Historie, er fundet, kan Religionen erkjendes i
sit eiendommelige Forhold til begge. Det Grundbegreb,
hvortil Vexelvirkningen af Mythologi, Historie og Religion
maa henføres, er Begrebet Frihed. Frihedens Princip
er Villie, Villiens Væsen er Aand; i Frihedens Verden
er det Aand og Villie, der behersker og bestemmer
den ydre Virkelighed. Er Folke-Aanden vaagen og Villien
stærk, saa er Handlekraften stor, og Historien rig
paa Bedrifter; er Aanden døsig og Villien mat, saa er
Handlekraften ringe, og den virkelige Historie ubetydelig.
Da nu Mythologien er det Aandens Speil, hvori
Folket seer sine Livsidealer og Formaal, saa angiver
denne Folke-Aandens første naturlige Selvaabenbarelse
tillige den Høide, hvortil et Folk ved egne Kræfter
er i Stand til at hæve sig, det Maal af Frihed, hvortil
dets virkelige Historie kan naae, og hvori den naturligen
maa have sin Grændse. Saaledes naaede Grækerne
i den gamle Verden netop til at virkeliggjøre det
Frihedsideal, som under Skjønhedens Form udtaler sig
i deres Mythologi; men heller ikke et Skridt længer.
Hvad der gjælder om Skjønhedsidealet hos Grækerne,
gjælder paa tilsvarende Maade om Magtidealet hos Romerne.
Udvortes seet, er det Barbarernes Vaaben, der
gjør Ende paa Romermagten; men, indvortes seet, havde
denne Magt allerede fra Begyndelsen af sit Maal, sin
Grændse i sig selv.

% --- p. 434 ---
\setcounter{footnote}{0}%
\opage{434}Dersom der nu i Virkeligheden ikke gaves anden
Aand end Folke-Aanden, saa vilde det see sørgeligt ud
med Friheden. Enhver Folkeaand har nemlig i Folkets
Individualitet en Naturgrændse, der ikke kan overskrides.
Vel er Aanden som Aand altid det Almene, men
dette Almene faaer kun Virkelighed med Kjød og Blod,
ved at leve og boe, i det Individuelle. At ville frigjøre
Aanden fra den naturlige Indskrænkning ved at opløse
den nationale Eiendommelighed i kosmopolitisk Ensformighed,
er, hvad Erfaringen noksom har viist, en metafysisk
Vildfarelse, et Forsøg paa at sætte Skyggen i
Legemets Sted og gjøre Livet selv til et blodløst Abstraktum.

Det ligger i Sagens Natur, at en tiltagende Civilisation
og dermed fremskridende Tænkning, efterhaanden
maa undergrave den mythiske Folketro og forvandle
Gudeskikkelserne til bevidste Symboler. Dersom nu en
Mythologi kunde tænkes at indeholde den høieste Sandhed
og angive Principet for den virkelige Frihed: da
vilde den Omdannelse, som Mytherne undergik i Reflexionen,
jo vistnok ikke betyde en Afkræftelse, men en Forhøielse
af Folke-Aanden, som paa denne Maade kom til
Klarhed over sig selv. Men saaledes forholder det sig
ikke. Det er ikke alene Gudernes udvortes Virkelighed,
Forstandskritiken angriber; den angriber Gudeverdenens
Væsen, dens Sandhed, idet Fornuften indseer, at et saadant
Væsen ikke længer fremmer, men hemmer Friheden.
Saa følger Opløsningen som af sig selv; thi den
bevidste Tænkning, der opdager det Ufornuftige og dermed
det Aandløse i Gudelæren, er ikke i Stand til at indsætte
en ny Aand istedenfor den gamle. Saaledes kunde
% --- p. 435 ---
\setcounter{footnote}{0}%
\opage{435}den græske Filosofi vel opløse den græske Mythologi,
men Opløsningen angreb Folke-Aanden; Filosofien selv
var magisk behersket af Folke-Aandens Grundformer:
ved at opløse Mythologien naaede den sin Fuldendelse,
men just som den stod i Begreb med at gaae videre,
indledede den sin egen Opløsning, thi Grændsen var sat
af Aanden.

Med den aabenbarede Religion kommer Sandhedens
Aand, den Villiens og Frihedens Aand, som alene formaaer
at løse Virkelighedens Gaader. Religionens Forhold
til Mythologien er fra dette Synspunkt af at bestemme
ved den aabenbarede Aands Forhold til Folke-Aanden.
For at see, hvad heri ligger, kunne vi begynde
med det enkelte Menneske. Den aabenbarede Religion
er hævet over al national Begrændsning; den er bestemt
for alle Folkeslag, altsaa for alle Mennesker, for
ethvert enkelt Menneske uden Undtagelse. Men uagtet
ethvert Menneske maa kunne tilegne sig Religionen, saa
er der dog, hvad Dybden, Inderligheden, Kraften og Selvstændigheden
i en saadan Tilegnelse angaaer, en umiskjendelig
Forskjel mellem Individerne. Individualiteten
er mere end et passivt Medium for Aandens Gjennemgang,
den er som Ædelstenen, der bryder Lyset og
kaster Straalerne tilbage med eiendommelig Glands og
Farvespil. Det kan uden Overdrivelse siges, at enhver
sand og dygtig Individualitet opfatter de religiøse Sandheder
paa eiendommelig Maade. Forskjellen hidrører
ikke ligefrem fra det Udvortes, thi Bekjendelsen kan
være den samme og Læren den samme; den væsenlige
Forskjel er i det Indvortes. Ikke blot i det Udvortes,
men ogsaa i det Indvortes er en Mand forskjel%
% --- p. 436 ---
\setcounter{footnote}{0}%
\opage{436}lig fra en Kvinde; hvoraf igjen følger, at det religiøse
Liv vinder en anden Skikkelse i Manden end i Kvinden.
Dersom den aabenbarede Religion ikke lagde den
allerstørste Vægt paa Bevarelse og Forevigelse af det i
Frihedslivet Individuelle, vilde Læren om de Dødes Opstandelse
være uden Betydning. --- Men Individualiteten i
den Enkelte er i det Hele uadskillelig fra Folke-Individualiteten.
At der saavel i denne som i hin er Meget,
der tilhører Endeligheden og hjemfalder til Forkrænkeligheden,
skal ikke nægtes --- det gjælder da især
om det Ydre --- men i det Almindeliges Navn maa Individualiteten
hævdes. Holdes nu Folke-Individualitetens
Betydning fast, da er der jo heller ikke noget Urimeligt
i, at eet Folk kan være anlagt paa en dybere Opfattelse
af Frihedslivet end et andet; den nordiske Mythologi
bærer i denne Henseende Vidne om de nordiske
Folks eiendommelige Begavelse. At gjøre et folkeligt
Fortrin gjældende, er ingenlunde det Samme som at
smigre Folkets Forfængelighed; det kan tverimod være
et Forsøg paa at vække dets Samvittighed. Det gjælder
her som overalt: »Af den, hvem Meget er betroet,
kan ogsaa meget fordres«. Naar en stor Fortid staaer
i tragisk Strid med en kummerlig Nutid, da er der Noget
iveien, der er Noget gaaet tabt. Historien vidner
imod det paagjældende Folk, at det har forsømt den
beleilige Tid og ikke brugt sine store Kræfter.

Jo dybere en Mythologi er, desto dybere er Folke-Aanden,
og jo dybere Folke-Aanden er, desto inderligere
og kraftigere vil Folket kunne tilegne sig Aabenbaringens
Aand. Det er Frihedslivet, altsaa Virkeligheden,
den indre Virkelighed, hvorpaa Alt først og sidst
% --- p. 437 ---
\setcounter{footnote}{0}%
\opage{437}kommer an. Den, der ikke kan forstaae, at de mythiske
Symboler endnu den Dag idag have Noget med det virkelige
Liv at gjøre, vil, naar vi see nærmere til, heller
ikke kunne føle og fatte, at Aabenbaringens overnaturlige
Kjendsgjerninger have Noget med Nutidens aandelige
Vilkaar at gjøre: Sandheden, den indre Virkelighed
i hine som i disse, er Aandsvirkelighed, Virkelighed
for Aanden. Men ligesom der er Forskjel paa timeligt
og evigt Liv, paa timelig og evig Frihed, paa Menneskets
Egenvillie og Guds almægtige, hellige Villie, saaledes
er der ogsaa Forskjel paa Folke-Aanderne og
Religionens hellige Aand; og da ogsaa en himmelvid
Forskjel mellem den Ide-Virkelighed, der kan tillægges
de mythiske Symboler, og den ideale Virkelighed, der
tilkommer det Overnaturlige, en principiel Forskjel mellem
Mythologi og Aabenbaring\footnote[1]{Om Religion og Mythologi jvfr. Forf.s Religionsfilosofi S. 11, flgd.}.

Den Grundtanke, som her er antydet i flygtige Omrids,
skulde nu i Folkeskolen antage Skikkelse, Liv og
Farve; dertil behøves Lærekræfter. Det koster allerede
Aandsanstrængelse og vedholdende Tænkning at arbeide
sig alsidig og grundig ind i en rig Verdensanskuelse, og
dog er dette endda kun den halve Vei til Maalet; thi
Vanskeligheden er at faae det Indvundne udfoldet og
fremstillet. Vi maae i denne Anledning gjøre opmærksom
paa den kategoriske Forskjel, der er imellem at
docere, at snakke gemytlig og at tale personlig.

At docere er at foredrage en Sag objektivt, klare
dens Sammenhæng i almene Begreber og vise dens Gyldighed
i Kraft af uimodsigelige Grunde. I det strenge

% --- p. 438 ---
\setcounter{footnote}{0}%
\opage{438}Kathederforedrag holder Læreren sin Subjektivitet tilbage;
hvis han i enkelte Yttringer lader den skinne igjennem,
da skeer det saa indirekt, at de forbigaaende Glimt af
individuel Stemning kun tjene til at forstærke Indtrykket
af Sagens Betydning og skærpe Opmærksomheden for
Grundenes Vægt. Kun ved at lægge Baand paa sin
Individualitet og fjerne uvedkommende Hensyn til private
Meninger kan Læreren tilveiebringe den ophøiede
Ro, den kjølige Besindighed, den Bevidsthedens ufordunklede
Klarhed, der er en Betingelse for, at Tilhøreren
uforstyrret kan følge Tankebevægelsen Punkt for
Punkt.

At snakke gemytlig er nu lige det Modsatte. Den
naturfriske, umiddelbare Subjektivitet kryber ikke i
Skjul bag det Objektive, men giver sig ganske som
den er, ubekymret, om det skeer i Søndagsdragt eller
in \textit{puris naturalibus} eller i Slaabrok og Tøfler. Det
gemytlige Foredrag kan spille i alle Farver, klimpre paa
alle Strenge og gjøre en høist forskjellig Virkning. Det
kan være tiltrækkende og underholdende, hvis den talende
eller, for at blive i Kategorien, snakkende Individualitet
har Lune, Vid og Humor i Overflødighed; men hvis disse
gode Gaver mangle, kan det ogsaa blive langtrukket,
besværligt og kjedeligt. Det kan blive til noget Belærende,
hvis den, der spøgende spiller med Ordet, har
Indsigt, Blik og Forstand, men i modsat Fald kan det
jo ogsaa blive til noget Andet, f. Ex. til Vrøvl. Over
alt dette raader en høiere, uransagelig Tilfældighed.

Den, der vil tale personlig, maa nu hverken holde
sig i Objektivitetens Kulde eller varme sig ved Blus af
gemytlige Indfalds Fyrværkeri. Han tør, hvad det Ob%
% --- p. 439 ---
\setcounter{footnote}{0}%
\opage{439}jektive angaaer, ikke byde sine Tilhørere en udenfra
kommende Belæring og lade dem forsøge at iføre
sig en Theori, der ikke særlig er for dem, men for
Menneskeheden, ligesom man lader Rekruten prøve en
Uniform, der ikke er syet til ham, men til Armeen;
han maa allerførst, om vi saa maae sige, have taget Maal
af sig selv og sine Tilhørere, saa at han virkelig er i
Stand til at sætte sig paa Tilhørernes Standpunkt, digte
sig ind i deres Forestillingskreds, tale ud af deres egen
Bevidsthed, røre Strengene i deres eget Bryst og hjælpe
dem til at tænke deres egne Tanker. Dette er Opgaven,
seet fra den subjektive Side. Men Sagen, det Objektive,
er ham ligesaa vigtig. Han maa i denne Henseende
medbringe en Fylde af Kundskaber, et rigt
Stof og have Stoffet tilbørlig gjennemarbeidet. Det
er ikke Emnet for sig, saalidt som det er Formen af
af blot individuel Eiendommelighed, der giver Talen
det personlige Præg: det er den aandige Enhed af
begge. Læreren maa i den Grad have optaget Sagen
i sin Bevidsthed, at han bogstavelig lever i den. Det
Hele maa være ham saa nærværende, at han i Talens
Øieblik, under Foredragets Gang, efter Tanke og
Stemning kan gribe Exemplerne frit, kan beherske og
forme det Mangfoldige saaledes, at et folkeligt Karakterbillede,
et Sjælemaleri af Historiens virkelige Liv,
uvilkaarlig danner sig for Tilhørerne. Saa bliver det
Høieste enfoldigt og forstaaes enfoldigt; Tanken iklædes
Fortællingen; Billeder og Reflexioner, Exempler og
Beskrivelser vexle med fuld Frihed, men uden Omsvøb
eller udsvævende Vidtløftighed: saa virker den hele
Fremstilling til at sikkre Forstaaelsen og fremme Er%
% --- p. 440 ---
\setcounter{footnote}{0}%
\opage{440}kjendelsens Udvikling, saa kan den Ulærde naturlig
følge med, og den Lærde have sin Glæde af at sidde
hos og lytte efter. Den bare Livlighed kan ikke gjøre
det, saalidt som den tørre Grundighed. Man er paa
Aandslærens Enemærker altfor tilbøielig til at sætte
Skilsmisse mellem det grundige og det vækkende Foredrag; idet man antager, at det grundige udelukkende
tilhører Videnskaben og derfor ikke kan nedlade sig til
at være vækkende, medens det vækkende, der falder
udenfor Videnskaben, umulig kan være grundigt, men
maa nøies med at være livligt. Det er en falsk Synsmaade.
Sandheden er, at et folkeligt Foredrag maa
holdes i en anden Form end et videnskabeligt; men de,
som antage, at et Foredrag er ugrundigt, fordi det er
vækkende og levende, maae vel ogsaa antage, at et
Foredrag altid er videnskabeligt og grundigt, naar det
blot er dygtig trægt, talentløst og kjedsommeligt.

Af det her Afhandlede vil det forhaabenlig nu være
klart, at Fremskridt, der gjøres i de nye Folkeskoler,
ikke beroe paa Elevernes Forkundskaber, men
paa Lærernes Duelighed. I Kundskabsskolen falder den
overveiende Del af Arbeidsbyrden paa Eleverne, thi det
er virkelig lettere at høre i en Lektie end at lære
en Lektie; i Folkeskolen derimod falder Arbeidsbyrden
fornemmelig paa Læreren, saasandt det altid er lettere
at høre paa et frugtbart vækkende Foredrag, end det er
at holde det. Kan det folkelige Foredrag blot udvikle sig
til den Klarhed og Fuldkommenhed, som Sagens Vigtighed
kræver, saa kan man være forvisset om, at Tilhørerne
nok skulle komme i Arbeide, at enhver mulig
Veiledning, den være mundlig eller skriftlig, den være
% --- p. 441 ---
\setcounter{footnote}{0}%
\opage{441}at finde i mindre Lærebøger eller i større Værker, skal   % PRINTER'S DEFECT (p. 441): „finde“: the f did not ink its head; letter supplied
efter Tid og Leilighed blive benyttet, thi driftig Selvvirksomhed
med friske Kræfter udretter store Ting.

Det har i nærværende Afhandling ikke været vor
Opgave enten at revidere givne Undervisningsplaner eller
gjøre Udkast til nye; vi have, med særligt Hensyn paa
Ide og Virkelighed, indskrænket os til at henlede Opmærksomheden
paa et stort Forehavende, som, takket være
de mange unge Mænd, der ikke alene hos os, men ogsaa
i Broderlandene med opoffrende Begeistring have lagt
Haand paa Værket, vil med Styrelsens Bistand bære
Frugter for aandelig Frihed og borgerligt Fremskridt
til Held og Velsignelse for Nordens Lande. Vi have
fæstet vort Blik paa den Vækkelse, paa hvilken der
efter vor Overbevisning alene kan bygges et Fremtids
Haab; men vi have ikke villet enten oversee eller miskjende
det Gode og Gavnlige, som Kundskabsskolen
saavel i lavere som høiere Stil har ydet og fremdeles
vil yde: vi have med den principielle Forskjel mellem
lærd og folkelig Dannelse for Øie villet bidrage til en
Erkjendelse af, at der i Nutiden kan og maa gives
to Slags Undervisning.
\signature{R. \emph{Nielsen}.}

\end{document}
